% Paper 32: The GeoVac Spectral Triple — A Synthesis
% GeoVac Project, May 2026
% v1.0 -- Synthesis paper. Constructs the spectral triple (A, H, D)
% explicitly on the Fock-projected S^3 graph, checks axioms, places
% the project inside the Marcolli-van Suijlekom gauge-network lineage,
% and reads Papers 25, 28, 30, 31 as structural sub-sectors.

\documentclass[11pt]{article}
\usepackage{amsmath,amssymb,amsthm}
\usepackage{booktabs}
\usepackage{array}
\usepackage[margin=1in]{geometry}

\newtheorem{theorem}{Theorem}
\newtheorem{proposition}{Proposition}
\newtheorem{corollary}{Corollary}
\newtheorem{observation}{Observation}
\newtheorem{definition}{Definition}
\newtheorem{remark}{Remark}

\title{The GeoVac Spectral Triple:\\
A Synthesis Paper Locking the Framework into the\\
Marcolli--van~Suijlekom Gauge-Network Lineage}
\author{GeoVac Project}
\date{May 2026}

\begin{document}
\maketitle

\begin{abstract}
The GeoVac framework has, over Papers~0--31, accumulated a body of
structural results---the $S^3$ Fock-projection equivalence
(Paper~7~\cite{paper7}), the angular sparsity theorem
(Paper~22~\cite{paper22}), the Hopf-graph lattice-gauge structure
(Paper~25~\cite{paper25}), QED on $S^3$ (Paper~28~\cite{paper28}),
the Bargmann--Segal $S^5$ HO lattice (Paper~24~\cite{paper24}), the
$\mathrm{SU}(2)$ Wilson gauge (Paper~30~\cite{paper30}), and the
universal/Coulomb partition (Paper~31~\cite{paper31})---each
self-contained but together suggesting one parent structure.  This
paper names and constructs that structure: a finite spectral triple
$(\mathcal{A}_{\mathrm{GV}}, \mathcal{H}_{\mathrm{GV}},
D_{\mathrm{GV}})$ on the Fock-projected $S^3$ graph, with
$\mathcal{A}_{\mathrm{GV}}$ the algebra of complex functions on graph
nodes, $\mathcal{H}_{\mathrm{GV}}$ the Camporesi--Higuchi spinor space
restricted to the truncation, and $D_{\mathrm{GV}}$ a graph Dirac
operator whose continuum limit is the Camporesi--Higuchi
operator~\cite{camporesi_higuchi1996} on unit~$S^3$.  This is not a
new derivation; it is a structural map.  Three results follow.
(1)~\emph{Lineage.}~The construction sits inside the published
Marcolli--van~Suijlekom gauge-network framework~\cite{marcolli_vs2014}
(with the Perez-Sanchez correction~\cite{perez_sanchez2024,%
perez_sanchez2025} that the continuum limit is Yang--Mills without
Higgs), and inside the Connes--van~Suijlekom spectral-truncation
program~\cite{connes_vs2021,hekkelman2022,hekkelman2024}.  Every
structural ingredient of the GeoVac construction has published
precedent; what is not duplicated is the $\alpha$ prediction at
$8.8 \times 10^{-8}$ (Paper~2~\cite{paper2}) from a discrete spectral
construction.
(2)~\emph{Sub-sector identification.}  Paper~25's Hodge-1 Laplacian is
the Cartan-subalgebra (maximal-torus) projection of an
almost-commutative extension of the present triple by~$M_2(\mathbb{C})$;
Paper~30's $\mathrm{SU}(2)$ Wilson action is the full non-abelian
slice of the same extension.  Paper~28's QED on~$S^3$ is the
spectral-action computation on $D_{\mathrm{GV}}$; Paper~31's
universal/Coulomb partition is the $\mathcal{A}/D$ split of the
present triple.  Papers~25, 28, 30, 31 are not four parallel
results---they are four projections of the single triple constructed
here.
(3)~\emph{Axiom audit.}  $\mathcal{A}_{\mathrm{GV}}$ is finite,
involutive, faithfully represented on $\mathcal{H}_{\mathrm{GV}}$;
$D_{\mathrm{GV}}$ is self-adjoint with bounded $[D_{\mathrm{GV}}, a]$
for $a \in \mathcal{A}_{\mathrm{GV}}$; the order-one condition is
trivially satisfied by the abelian core; a real structure
$J_{\mathrm{GV}}$ exists on the Dirac sector with
$J_{\mathrm{GV}}^2 = -\mathbb{1}$ (the quaternionic/pseudoreal sign
for the 2-dimensional Dirac spinor representation of
$\mathrm{Spin}(3) = \mathrm{SU}(2)$) and $J_{\mathrm{GV}} D_{\mathrm{GV}}
= +D_{\mathrm{GV}} J_{\mathrm{GV}}$, matching the standard signs
$(\epsilon, \epsilon') = (-1, +1)$ for $\mathrm{KO}$-dimension~$3$
mod~$8$ of~$S^3$~\cite{connes1995}; a $\mathbb{Z}_2$-grading
$\gamma_{\mathrm{GV}}$ is absent at the Dirac level (correctly:
$S^3$ is odd-dimensional, the spectral triple is \emph{odd}).  The four-way $S^3$ coincidence---Fock projection image,
Hopf bundle base, Camporesi--Higuchi spin carrier,
$\mathrm{SU}(2)$ gauge manifold---is the statement that the manifold
of $\mathcal{A}_{\mathrm{GV}}$ in the continuum limit, the spinor
bundle of $D_{\mathrm{GV}}$, and the gauge group of natural
fluctuations of $D_{\mathrm{GV}}$ all coincide on the unique rank-1
non-abelian compact simply-connected Lie group.

\noindent\textbf{Scope.}  This paper does not derive new physics.  It
synthesizes existing GeoVac results into the axiomatic spectral-triple
language.  The synthesis lifts each result from a free-standing
observation to a sub-sector of a single named structure, and makes
explicit which axioms hold rigorously, which hold up to the standard
finite-truncation caveats, and which remain open---specifically
whether the GeoVac Dirac graph operator carries a real structure
matching the topological $\mathrm{KO}$-dimension of $S^3$ exactly at
finite $n_{\max}$ or only in the continuum limit.
\end{abstract}

% ===================================================================
\section{Introduction}
\label{sec:intro}
% ===================================================================

\subsection{Why a spectral triple, and why now}

Across thirty-one papers the GeoVac framework has carried four
distinct strands.  Strand~1 is the Fock projection
$p_0 = \sqrt{-2E_n}$ that maps the hydrogenic Hamiltonian onto the
free Laplacian on the unit three-sphere, with eigenvalues
$\lambda_n = -(n^2 - 1)$ on graph nodes labeled by quantum-number
triples $(n, l, m_l)$ (Paper~7).  Strand~2 is the Hopf bundle
$S^1 \to S^3 \to S^2$ that organizes the discrete edge structure into
ladder operators with $U(1)$ phase content (Papers~2, 25).  Strand~3
is the Camporesi--Higuchi Dirac spectrum
$|\lambda_n^{\mathrm{D}}| = n + 3/2$ with degeneracy
$g_n^{\mathrm{D}} = 2(n+1)(n+2)$ that controls QED on $S^3$
(Paper~28).  Strand~4 is the Wilson lattice gauge structure with
$\mathrm{SU}(2)$ link variables on the same graph (Paper~30).

Each strand uses the unit three-sphere in a different role:
\begin{itemize}\setlength{\itemsep}{0pt}
\item[(a)] $S^3$ as the \emph{conformal image} of the hydrogenic
spectrum (Paper~7);
\item[(b)] $S^3$ as the \emph{base} of the bundle carrying the
fine-structure-constant construction (Paper~2);
\item[(c)] $S^3$ as the \emph{spin carrier} of the Camporesi--Higuchi
Dirac operator (Paper~28);
\item[(d)] $S^3$ as the \emph{gauge manifold} of $\mathrm{SU}(2)$
Wilson links (Paper~30).
\end{itemize}
Working hypothesis~WH4 in the project register flags this as the
strangest and most consequential coincidence in the framework: $S^3 =
\mathrm{SU}(2)$ is the unique rank-1 non-abelian compact simply
connected Lie group where (a)--(d) all coincide; in any other
dimension or any other rank, the four roles separate.  WH4 takes the
position that the four strands are not four coincidences but a single
spectral-triple structure viewed in four projections.

This paper makes WH4 explicit.  We give a concrete construction of a
finite spectral triple
$(\mathcal{A}_{\mathrm{GV}}, \mathcal{H}_{\mathrm{GV}},
D_{\mathrm{GV}})$ (Sec.~\ref{sec:construction}); audit the standard
Connes axioms one by one (Sec.~\ref{sec:axiom_audit}); identify
Papers~25, 28, 30, and~31 as structural sub-sectors of the resulting
object (Sec.~\ref{sec:subsectors}); read the Paper~2 combination
formula $K = \pi(B + F - \Delta)$ inside the spectral-action language
(Sec.~\ref{sec:K_decomposition}); and place the Coulomb/HO asymmetry
of Paper~31 in spectral-triple terms (Sec.~\ref{sec:coulomb_ho}).

\subsection{What this paper is and is not}

\paragraph{What it is.}  A synthesis paper.  The construction is
explicit, finite, and reproducible from existing GeoVac modules
(\texttt{geovac/lattice.py}, \texttt{geovac/dirac\_matrix\_elements.py},
\texttt{geovac/fock\_graph\_hodge.py}); the lineage citations are to
published mathematical work; the audit of the Connes axioms uses
standard textbook definitions~\cite{connes1995,connes_marcolli2008,%
suijlekom_book2015}.

\paragraph{What it is not.}  A derivation of $\alpha$.  Paper~2 stays
where it is: a structural observation with three independent spectral
homes for the three pieces $B$, $F$, $\Delta$ but no first-principles
derivation of the combination rule.  This paper does not change that
status; it gives the spectral-action language in which the open
question is phrased without changing the conjectural label.

\paragraph{What it is not, second.}  A claim of the
Connes--Chamseddine Standard Model.  The GeoVac construction is a
\emph{rank-1} non-abelian almost-commutative extension; the SM
spectral triple of \cite{chamseddine_connes2010} is rank~$\ge 3$ with
specific finite-algebra content $\mathbb{C} \oplus \mathbb{H} \oplus
M_3(\mathbb{C})$.  GeoVac is not the Standard Model.  It is a finite
spectral triple in the same lineage, with the empirically interesting
property that one of its spectral-action coefficient combinations
matches~$\alpha^{-1}$ at $8.8 \times 10^{-8}$.

\paragraph{Internal register.}  This paper makes explicit the
spectral-triple framing already operating informally in the working
hypothesis register (\texttt{CLAUDE.md}~\S 1.7, WH1 ``MODERATE--STRONG'').
Papers continue under the rhetoric rule of \S 1.5: dual-description
framing, no ontological priority claims, lead with computational
results.

% ===================================================================
\section{The Connes Data Recalled}
\label{sec:connes_data}
% ===================================================================

A spectral triple is a triple $(\mathcal{A}, \mathcal{H}, D)$
where~\cite{connes1995}
\begin{itemize}\setlength{\itemsep}{2pt}
\item $\mathcal{A}$ is an involutive (unital) algebra,
\item $\mathcal{H}$ is a Hilbert space carrying a faithful
$*$-representation $\pi: \mathcal{A} \to \mathcal{B}(\mathcal{H})$,
\item $D : \mathcal{H} \to \mathcal{H}$ is an unbounded self-adjoint
operator with compact resolvent, such that the commutators
$[D, \pi(a)]$ are bounded for every $a \in \mathcal{A}$.
\end{itemize}
The triple is called \emph{even} if there is a self-adjoint involution
$\gamma$ commuting with $\pi(\mathcal{A})$ and anticommuting with $D$;
otherwise \emph{odd}.  A \emph{real structure} on the triple is an
antilinear isometry $J: \mathcal{H} \to \mathcal{H}$ with
$J^2 = \epsilon\,\mathbb{1}$, $JD = \epsilon'\,DJ$,
$J\gamma = \epsilon''\,\gamma J$ (the signs $(\epsilon,\epsilon',
\epsilon'')$ are determined by the $\mathrm{KO}$-dimension
mod~$8$~\cite{connes1995}); for the canonical Riemannian spin
manifold of dimension $n$, the $\mathrm{KO}$-dimension is $n$ mod~$8$
and the signs are tabulated.

The two structural axioms of central interest below are:
\begin{itemize}
\item[(\textbf{O1})] \emph{Order one:} $[[D, \pi(a)], \pi(b)^{\circ}]
= 0$ for all $a, b \in \mathcal{A}$, where $b^{\circ} = JbJ^{-1}$ is
the right action induced by the real structure;
\item[(\textbf{R})] \emph{Reality:} $J$ exists with the signs above.
\end{itemize}

For an \emph{almost-commutative} spectral triple
$(\mathcal{A}_M \otimes \mathcal{A}_F, \mathcal{H}_M \otimes
\mathcal{H}_F, D_M \otimes \mathbb{1} + \gamma_M \otimes
D_F)$~\cite{vansuijlekom_book2015,chamseddine_connes2010}, the
manifold sector $\mathcal{A}_M = C^{\infty}(M)$ is replaced in our
finite setting by the algebra of complex-valued functions on a finite
graph that converges to $C^{\infty}(M)$ in the continuum limit.  This
is exactly the move made by Marcolli and van~Suijlekom in their gauge
networks construction~\cite{marcolli_vs2014} and continued by
Perez-Sanchez~\cite{perez_sanchez2024,perez_sanchez2025}: replace the
manifold algebra by an algebra on a discrete vertex set, with edges
carrying connection data.  The GeoVac construction sits in this
lineage by design.

% ===================================================================
\section{The GeoVac Triple: Explicit Construction}
\label{sec:construction}
% ===================================================================

Fix a truncation $n_{\max} \ge 1$.  All objects below are finite.

\subsection{The algebra $\mathcal{A}_{\mathrm{GV}}$}

Let $V_{\mathrm{Fock}}$ denote the set of Fock-graph nodes at
truncation $n_{\max}$, i.e. quantum-number triples $(n, l, m_l)$ with
$1 \le n \le n_{\max}$, $0 \le l \le n-1$, $-l \le m_l \le l$.  The
total node count is $N_{\mathrm{Fock}}(n_{\max}) = \sum_{n=1}^{n_{\max}}
n^2$.

\begin{definition}[Function algebra on the Fock graph]
\label{def:A_GV}
$\mathcal{A}_{\mathrm{GV}}^{(n_{\max})} := \mathbb{C}^{V_{\mathrm{Fock}}}$,
the commutative $*$-algebra of complex-valued functions on
$V_{\mathrm{Fock}}$ under pointwise multiplication, with involution
$f^* (v) = \overline{f(v)}$.
\end{definition}

This is precisely the manifold algebra of the gauge networks of
\cite{marcolli_vs2014} (their $C(V)$, with $V$ the vertex set), with
the GeoVac-specific choice that $V$ is the Fock-projected $S^3$ graph
of Paper~7 rather than a generic graph.  It is finite-dimensional
($\dim_{\mathbb{C}} \mathcal{A}_{\mathrm{GV}} = N_{\mathrm{Fock}}$),
unital ($1$~is the constant function~$1$), and abelian.

\begin{remark}[$\mathcal{A}_{\mathrm{GV}}$ as $C^*$-envelope of an operator system]
\label{rem:operator_system}
In the parallel reading of GeoVac through the spectral-truncation
framework of Connes and van~Suijlekom~\cite{connes_vs2021},
Definition~\ref{def:A_GV} describes
$\mathcal{A}_{\mathrm{GV}}^{(n_{\max})} = \mathbb{C}^{V_{\mathrm{Fock}}}$
as the $C^*$-envelope $C^*(\mathcal{O}_{n_{\max}})$ of the operator
system
\[
   \mathcal{O}_{n_{\max}} := P_{n_{\max}} \, C^\infty(S^3) \, P_{n_{\max}},
\]
where $P_{n_{\max}}$ is the index-truncation projecting onto the
span of the first $n_{\max}$ Fock shells (in the spinor-harmonic
basis of $L^2(S^3)$ adapted to the SO(4) decomposition).  This is
structurally identical to the Connes--van~Suijlekom worked example
on $S^1$ (their $P_n$ onto $\mathrm{span}_{\mathbb{C}}\{e_1, e_2,
\ldots, e_n\}$, \cite[\S 3.1]{connes_vs2021}): a basis-index cutoff,
not a Borel functional-calculus projection
$\chi_{[-\Lambda, \Lambda]}(D)$.  The genuine truncated algebra-image
$\mathcal{O}_{n_{\max}}$ is $*$-closed but \emph{not} multiplicatively
closed; it carries nontrivial off-diagonal matrix elements
\[
   \langle (n,l,m_l) | \, f \, | (n', l', m_l') \rangle
   = \int_{S^3} \overline{Y_{n l m_l}}\, f \, Y_{n' l' m_l'}\, dV
\]
between distinct $n$-shells.  These overlap integrals are exactly
the Gaunt~/~Wigner-3j angular machinery and the hypergeometric
Slater radial integrals already implemented throughout the GeoVac
pipeline (Paper~22 angular sparsity; Paper~14 ERI structure).  The
off-diagonal data therefore already exists in the framework;
$\mathcal{A}_{\mathrm{GV}}$ as written in Definition~\ref{def:A_GV}
is its diagonal restriction.

The two readings agree on diagonal elements and on the spectral
data of $D_{\mathrm{GV}}$; they differ in what is taken as the
algebraic content surviving the truncation.  Connes and
van~Suijlekom note that taking the $C^*$-envelope of $\mathcal{O}_\Lambda$
typically collapses to a ``non-informative full matrix algebra'' --
the algebraic-structure information lives in the operator-system
data thrown away by the envelope construction.  Constructing
$\mathcal{O}_{n_{\max}}$ explicitly in $\mathcal{B}(\mathcal{H}_{\mathrm{GV}})$
and computing its propagation number and Connes distance on
node-evaluation states is the deliverable of the WH1 round-2 sprint
(see \texttt{debug/wh1\_connes\_vs\_mapping\_memo.md}, 2026-05-03).
The present subsection retains the
$\mathcal{A}_{\mathrm{GV}} = \mathbb{C}^{V_{\mathrm{Fock}}}$ definition
because (a) it matches the Marcolli--van~Suijlekom gauge-network
convention used throughout Papers~25 and~30, and (b) the spectral
triple's structural verifications (Theorem~\ref{thm:GV_triple}
below; KO-dimension audit; Hodge structure) are unaffected by
which representative of the equivalence class is taken.
\end{remark}

\begin{proposition}[Propagation number of $\mathcal{O}_{n_{\max}}$]
\label{prop:propagation_number}
Let $\mathcal{O}_{n_{\max}} = P_{n_{\max}} C^\infty(S^3) P_{n_{\max}}$
be the truncated operator system of
Remark~\ref{rem:operator_system}, viewed as a $*$-closed subspace of
$M_N(\mathbb{C})$ with $N = N_{\mathrm{Fock}}(n_{\max})$.  The
propagation number in the sense of Connes \& van~Suijlekom
(\cite{connes_vs2021}, Definition~2.39) satisfies
\begin{equation}
   \mathrm{prop}(\mathcal{O}_{n_{\max}}) = 2
   \qquad \text{for every } n_{\max} \in \{2, 3, 4\},
   \label{eq:prop_two}
\end{equation}
matching the Toeplitz operator system $C(S^1)^{(n)}$
(\cite{connes_vs2021}, Proposition~4.2: $\mathrm{prop}(C(S^1)^{(n)}) =
2$ independently of $n$) verbatim.  Equivalently,
$\dim_{\mathbb{C}}(\mathcal{O}_{n_{\max}}^2) = N^2$ while
$\dim_{\mathbb{C}}(\mathcal{O}_{n_{\max}}) < N^2$ strictly:
\[
\renewcommand{\arraystretch}{1.15}
\begin{array}{c|c|c|c|c}
n_{\max} & N & N^2 & \dim(\mathcal{O}_{n_{\max}}) &
   \dim(\mathcal{O}_{n_{\max}}^2) \\ \hline
2 & 5  & 25  & 14  & 25  \\
3 & 14 & 196 & 55  & 196 \\
4 & 30 & 900 & 140 & 900 \\
\end{array}
\]
The dimension fraction $\dim(\mathcal{O}_{n_{\max}}) / N^2$ falls
$0.560 \to 0.281 \to 0.156$ across the three truncations: the
operator system is a vanishing fraction of its $C^*$-envelope as
$n_{\max}$ grows, but its second power fills the envelope
completely.
\end{proposition}

\begin{proof}[Computational verification]
The full construction and verification is documented in
\texttt{geovac/operator\_system.py} (the \texttt{TruncatedOperatorSystem}
class) and \texttt{tests/test\_operator\_system.py} (24 tests, all
passing; CPU time ${\sim}20$~s).  In summary: a basis of
$\mathcal{O}_{n_{\max}}$ is built from the multiplier matrices
$M_{NLM}$ corresponding to the hyperspherical harmonics
$Y^{(3)}_{NLM}$, with selection rules $|n - n'| + 1 \le N \le n + n' -
1$ (SO(4) triangle), $|l - l'| \le L \le l + l'$ with $l + l' + L$
even (S$^2$ Gaunt parity), $M = m - m'$, and the angular S$^2$ part
evaluated symbolically via \texttt{sympy.physics.wigner.wigner\_3j}.
Linear independence of $\{M_{NLM}\}$ is verified at machine precision;
the propagation number is read off as the smallest $k$ with
$\dim(\mathcal{O}_{n_{\max}}^k) = N^2$.  The result
$\mathrm{prop} = 2$ is robust: it is invariant under the choice of
SO(4) radial-overlap placeholder (verified across three placeholders
including a pure indicator that retains only the support pattern), so
\eqref{eq:prop_two} is a structural invariant of the SO(4) selection
rules at finite cutoff and does not depend on the radial value
convention.  See \texttt{debug/wh1\_round2\_propagation\_memo.md}
(2026-05-03) for the full computational record, including the explicit
witness pair $a = M_{2,1,0}$ at $n_{\max} = 2$ for which $a^2 \notin
\mathcal{O}_{n_{\max}}$ (Frobenius residual $\approx 14.9\%$), with
the deficit concentrated on the top-shell diagonal entries.
\end{proof}

\begin{remark}[Significance of $\mathrm{prop} = 2$]
Connes \& van~Suijlekom prove $\mathrm{prop}(C(S^1)^{(n)}) = 2$ for
the Toeplitz operator system on the circle by showing that products
of two Toeplitz shift generators reach every matrix unit
$e_{ij}$~\cite{connes_vs2021}.  The companion dual operator system
$C(\mathbb{Z})^{(n)}$ has $\mathrm{prop} = \infty$ (their
Proposition~4.3) — so the value $\mathrm{prop} = 2$ is non-trivial
even at this level: it is specific to the Toeplitz construction, not
to operator systems in general.  Eq.~\eqref{eq:prop_two} therefore
states that the Fock-projected $S^3$ truncation is in the same
propagation class as the Toeplitz $S^1$ — one of the simplest worked
examples in the spectral-truncation literature — at the level of a
quantitative Connes--van~Suijlekom invariant, not merely at the
level of the underlying construction.  Together with the
$\mathrm{SO}(4)$ isometry preservation discussed in
Sec.~\ref{sec:axiom_audit} below, this is the strongest
quantitative alignment between GeoVac and the published Connes-vS
framework currently known.

The mechanism producing
$\mathcal{O}_{n_{\max}} \subsetneq M_N(\mathbb{C})$ is the
``raise-then-lower at the top shell escapes the truncation'' boundary
correction: a multiplier $M_{N \ge 2, L, M}$ acting at the top shell
$n = n_{\max}$ would, in the continuum, be balanced by intermediate
sums over $n' > n_{\max}$, but the truncation $P_{n_{\max}}$ kills
these contributions.  This is structurally the same boundary effect
as the pendant-edge theorem of Paper~28~\cite{paper28},
$\Sigma_{\mathrm{graph}}(\mathrm{GS}) = 2(n_{\max} - 1)/n_{\max}$ for
the graph-native ground-state self-energy, viewed from the
operator-system side rather than the QED side.
\end{remark}

\begin{corollary}[Structural specificity of $\mathrm{prop} = 2$ on $S^3$]
\label{cor:structural_specificity}
The propagation value
$\mathrm{prop}(\mathcal{O}_{n_{\max}}) = 2$ in
Proposition~\ref{prop:propagation_number} is structurally specific to
the Connes--van~Suijlekom spectral truncation of $S^3$ and \emph{not}
a generic feature of arbitrary finite-dimensional truncations of $S^3$.
The natural circulant-style comparator on $S^3$ at $N$ points
\[
   A_{\mathrm{circ}, N} := \mathrm{span}\{E_{i, i} : i = 0, \ldots, N - 1\}
   \subset M_N(\mathbb{C}),
\]
i.e., the commutative $C^*$-subalgebra of diagonal matrices in the
``point-evaluation'' basis of $\mathbb{C}^N$, satisfies
\begin{equation}
   \mathrm{prop}(A_{\mathrm{circ}, N}) = \begin{cases}
      1 & \text{(intrinsic envelope:\ $A_{\mathrm{circ}, N}$ is its own
                  $C^*$-algebra)}, \\
      \infty & \text{(ambient envelope:\ $M_N(\mathbb{C})$ is non-abelian;
                       no power of an abelian subalgebra reaches it)}.
   \end{cases}
   \label{eq:prop_circulant}
\end{equation}
Verified at $N \in \{2, 5, 14, 30, 55, 140\}$, matching both
dimension-matching conventions (envelope-match $N(\mathrm{circ}) =
N(\mathrm{GV})$ and operator-system-match $N(\mathrm{circ}) =
\dim(\mathcal{O}_{n_{\max}})$).  Either reading of
\eqref{eq:prop_circulant} is categorically different from the
spectral-truncation value~$2$.
\end{corollary}

\begin{proof}[Computational verification]
The construction lives in \texttt{geovac/circulant\_s3.py} (the
\texttt{CirculantS3Truncation} class) with verification in
\texttt{tests/test\_circulant\_s3.py} (35 tests, all pass; runtime
${\sim}1.4$~s).  Algorithmically, $\mathrm{prop}(A_{\mathrm{circ}, N})$
is computed via the same vec-stack rank algorithm used in the
proof of Proposition~\ref{prop:propagation_number}, applied to the
diagonal-matrix basis of $A_{\mathrm{circ}, N}$.  The intrinsic-envelope
reading exits at $k = 1$ since $\dim(A_{\mathrm{circ}, N}^1) = N =
\dim_{\mathbb{C}}(A_{\mathrm{circ}, N})$ already saturates the
intrinsic $C^*$-envelope.  In the ambient-envelope reading, the
$M_N(\mathbb{C})$ envelope has dimension $N^2$ for $N \ge 2$ but
$\dim(A_{\mathrm{circ}, N}^k) = N$ for every $k$ (commutativity is
preserved under self-products), so no finite $k$ achieves
saturation; $\mathrm{prop} = \infty$ by Connes--van~Suijlekom
Definition~2.39.  See
\texttt{debug/wh1\_round3\_falsification\_memo.md} (2026-05-03) for
the full record, including the verbatim Connes--van~Suijlekom Outlook~\S6
contrast (lines 2611--2618 of arXiv:2004.14115v2) between
spectral and circulant truncations on~$S^1$, of which
Eq.~\eqref{eq:prop_circulant} is the verbatim $S^3$ analog.
\end{proof}

\begin{remark}[Connes distance on $\mathcal{O}_{n_{\max}}$: SDP framework
and structural findings]
\label{rem:connes_distance}
The Connes distance on operator systems
(\cite{connes_vs2021}, Eq.~(2))
\[
   d_D(\phi, \psi) = \sup_{x \in \mathcal{O}} \big\{ |\phi(x) - \psi(x)|
      : \|[D, x]\|_{\mathrm{op}} \le 1 \big\}
\]
restricted to pure node-evaluation states $\phi_v(x) := \langle v | x | v
\rangle$ on the truncated operator system has been computed by
semidefinite programming on $\mathcal{O}_{n_{\max}}$ at $n_{\max} = 2, 3$.
The implementation lives in \texttt{geovac/connes\_distance.py} (a
\texttt{cvxpy} SDP with the operator-norm constraint lifted to the LMI
$\bigl[\begin{smallmatrix} I & [D, x] \\ [D, x]^* & I \end{smallmatrix}\bigr]
\succeq 0$, and a gauge-fixing constraint orthogonalizing $x$ to
$\ker([D, \cdot]) \cap \mathcal{O}_h$); 17 tests in
\texttt{tests/test\_connes\_distance.py} (all pass; runtime ${\sim}9$~s).
Three structural findings emerged at the placeholder convention:

\begin{enumerate}
\item \textbf{SO(3) $m$-reflection forces zero distance.}  For every
   $|n, l, -m\rangle$ and $|n, l, +m\rangle$ pair, $d_D(\phi_{|n,l,-m\rangle},
   \phi_{|n,l,+m\rangle}) = 0$ exactly (one such pair at $n_{\max} = 2$,
   four at $n_{\max} = 3$).  This is the SO(3) $m$-reflection symmetry
   of the operator system showing through, not a pathology: the symmetry
   identifies the two states as members of the same orbit, and any
   point-discriminating function must be constant on the orbit.

\item \textbf{Clean rational structure at $n_{\max} = 2$.}  The five
   non-zero entries of the $5 \times 5$ distance matrix at $n_{\max} = 2$
   group into ratios $1 : 2 : 3$ (specifically $28.87 : 57.73 : 86.60$),
   with $86.6020 = 50\sqrt{3}$ to five decimals.  This rational/algebraic
   structure is lost at $n_{\max} = 3$ as the SDP gains degrees of
   freedom.  We log this as an algebraic fingerprint of the truncated
   metric at small cutoff; whether it persists in any natural
   generalization is open.

\item \textbf{Non-monotonicity in the Fock-graph distance.}  Pearson
   correlation between $d_D(\phi_v, \phi_w)$ and the combinatorial
   Fock-graph distance $d_{\mathrm{graph}}(v, w)$ (counting $\Delta n =
   \pm 1$ and $\Delta l = \pm 1$ ladder steps) is $\rho \approx -0.04$
   at $n_{\max} = 2$, growing to $+0.14$ at $n_{\max} = 3$ (Spearman
   $0.09 \to 0.18$).  The trend is weakly toward correlation but the
   absolute level is far from a discretized-geodesic metric at our
   reachable cutoffs.
\end{enumerate}

These findings should be read with three explicit caveats, none of
which are addressed in the present paper.

\begin{itemize}
\item The SO(4) radial overlap integral $I_{n N n'}^{l L l'}$ in the
   construction of $\mathcal{O}_{n_{\max}}$
   (Remark~\ref{rem:operator_system}) is currently a placeholder rather
   than the genuine Avery--Wen--Avery $3$-$Y$ integral on $S^3$.
   Round~2's propagation-number result
   (Proposition~\ref{prop:propagation_number}) is invariant under this
   choice; the Connes distance is not.  The absolute distance values
   reported above are convention-dependent until the placeholder is
   replaced.

\item The construction is on the scalar Fock basis $\mathcal{H}_{\mathrm{GV}}$
   only; the spinor lift to $\mathcal{H}_{\mathrm{GV}} \otimes
   \mathbb{C}^4$ with the native Camporesi--Higuchi Dirac in the
   $(n_{\mathrm{D}}, \kappa, m_j)$ basis is not yet incorporated.  The
   diagonal scalar Dirac on $\mathcal{H}_{\mathrm{GV}}$ has a large
   $\ker([D, \cdot]) \cap \mathcal{O}$ that makes the SDP unbounded on
   most state pairs; the metric values reported here use an
   off-diagonal Dirac proxy as a finite-distance regulator, which is a
   second convention choice.

\item No comparison to the round-$S^3$ continuum geodesic distance is
   attempted.  Such a comparison is the actual content of Connes--vS
   Gromov--Hausdorff convergence
   (Sec.~\ref{sec:axiom_audit} discussion of Row~9 of the round-1
   mapping), which has not been proved on $S^3$ in the published
   literature.
\end{itemize}

The structural findings (SO(3) $m$-reflection forced zeros; the $1 : 2
: 3$ rational ratios at $n_{\max} = 2$) are robust across the
placeholder choice; the absolute distance values and the
non-monotonicity diagnostic are not.  See
\texttt{debug/wh1\_round3\_connes\_distance\_memo.md} for the full
record.  Caveats (1) and (2) above are subsequently closed in
Remark~\ref{rem:r31_r32_update}; one of the three structural
findings is retracted there as a placeholder artifact.
\end{remark}

\begin{remark}[R3.1+R3.2 closure of round-3 caveats]
\label{rem:r31_r32_update}
Two of the three caveats above have been closed in subsequent
sub-sprints (continuation of the same dispatch on 2026-05-03), and
one of the three structural findings has been retracted as a
placeholder artifact.

\paragraph{R3.1 (Avery--Wen--Avery integral, placeholder $\to$
physical).}  The placeholder for $I_{n N n'}^{l L l'}$ is replaced
with the closed-form Avery--Wen--Avery $3$-$Y$ integral on $S^3$,
factorized as $I_{\mathrm{rad}} \cdot G_{\mathrm{Gaunt}}$ with
$I_{\mathrm{rad}}$ the Gegenbauer triple integral on $\chi$ and
$G_{\mathrm{Gaunt}}$ the standard $S^2$ Gaunt integral.
Implemented in \texttt{geovac/so4\_three\_y\_integral.py}
(sympy-exact, \texttt{lru\_cache}), verified by hyperspherical
orthonormality at $n_{\max} \le 3$ in exact arithmetic.
$\mathrm{prop}(\mathcal{O}_{n_{\max}}) = 2$ holds at $n_{\max} =
2, 3, 4$ under the physical default with dim sequences $14 \to
25,\ 55 \to 196,\ 140 \to 900$ \emph{bit-identical} to the
placeholder values---confirming the round-2 invariance claim by
replacement, not by sampling.

The Connes distance recomputed under physical values retains
the $m$-reflection forced zeros and the $50\sqrt{3}$ fingerprint,
but \emph{retracts} the $1 : 2 : 3$ rational structure as a
placeholder artifact: under Avery the distinct nonzero values at
$n_{\max} = 2$ are $\{1.27, 86.60, 87.18, 174.93, 175.53,
262.11\}$ with no clean integer grouping.  The $M_{2,0,0}$
multiplier rescales by $\sim 3\times$ between placeholder and
physical values, propagating as $\sim 9\times$ on pairs that use
it twice.  More importantly, the non-monotonicity \emph{flips
sign} at $n_{\max} = 3$: Pearson goes from $+0.14$ (placeholder)
to $-0.36$ (physical), Spearman from $+0.18$ to $-0.33$.  The
optimistic ``slow trend toward correlation'' reading is
retracted; under physical values the metric is
\emph{anti-correlated} with the Fock-graph distance.  Memo:
\texttt{debug/wh1\_r31\_avery\_wen\_avery\_memo.md}.

\paragraph{R3.2 (spinor lift, scalar Fock $\to$ Weyl spinor on
$S^3$).}  The construction extends to the Weyl sector of the
Camporesi--Higuchi spinor bundle on $S^3$ via Clebsch--Gordan
decomposition $|\psi_{l, j = l + 1/2, m_j}\rangle = \alpha_+
|Y_{l, m_j - 1/2}\rangle \otimes |\!\!\uparrow\rangle + \alpha_-
|Y_{l, m_j + 1/2}\rangle \otimes |\!\!\downarrow\rangle$ with
$\alpha_\pm = \sqrt{(l \pm m_j + 1/2)/(2l + 1)}$.  Each scalar
Avery multiplier $M_{NLM}$ lifts to a CG-weighted sum of two
scalar matrix elements; the resulting
$\widetilde{\mathcal{O}}_{n_{\max}}$ is API-compatible with the
existing SDP machinery.  Implemented in
\texttt{geovac/spinor\_operator\_system.py}.  $\dim
\widetilde{\mathcal{O}} = \dim \mathcal{O}$ at every $n_{\max}$
(the CG decomposition introduces no new linear independencies);
identity, $\ast$-closure, and $m_j$-reflection forced zeros all
generalize naturally.

Two clean R3.2 findings.  \textbf{(i)} \emph{The truthful
Camporesi--Higuchi Dirac on the spinor bundle is too
$n$-degenerate} for the Connes distance to be well-defined on
cross-shell pure-state pairs: at $n_{\max} = 2$, $24$ of $28$
cross-pairs give $+\infty$ because the $n$-shell-diagonal
multipliers $M_{N, L = \mathrm{even}, M = 0}$ commute with $D =
\mathrm{diag}(n_{\mathrm{Fock}} + 1/2)$, making the SDP
unbounded.  This is the spinor analog of the round-3
\texttt{shell\_scalar} diagnostic and confirms the
$n$-degeneracy obstruction is structural, not scalar-only.
\textbf{(ii)} \emph{Under a CH+offdiag perturbation that lifts
the $n$-degeneracy, the metric remains anti-correlated} with the
Fock-graph distance: Pearson nz $-0.36$ at $n_{\max} = 2$,
$-0.26$ at $n_{\max} = 3$.  The non-monotonicity is therefore
robust under both the placeholder $\to$ Avery upgrade (R3.1)
\emph{and} the scalar $\to$ spinor lift (R3.2).  Memo:
\texttt{debug/wh1\_r32\_spinor\_lift\_memo.md}.

\paragraph{Disambiguation settled.}  The round-3 \S 5.4
two-reading question (``placeholder artifact'' vs
``pure-state-distance pathology'') lands on the latter.  R3.1
rules out the placeholder reading (anti-correlation
\emph{strengthens} under physical values).  R3.2 rules out the
scalar-Dirac-proxy reading (anti-correlation persists under the
spinor lift).  The Connes distance on $\mathcal{O}_{n_{\max}}$ or
its spinor lift, with pure node-evaluation states, is structurally
not a discretization of the round-$S^3$ geodesic distance at any
cutoff we can reach.  The third caveat above (no GH comparison
on the full state space) remains the natural next move; whether
the metric becomes monotone in the GH limit via a Kantorovich--
Wasserstein lift to mixed states requires Peter--Weyl on
$\mathrm{SU}(2)$ (genuinely new NCG mathematics; deferred by
Connes \& van~Suijlekom 2021 three times, no published treatment
on $S^3$).
\end{remark}

\begin{remark}[Two-sided structural alignment]
\label{rem:two_sided_alignment}
Connes \& van~Suijlekom's worked $S^1$ example exhibits the
spectral-vs-circulant distinction along two axes simultaneously:
the Toeplitz operator system $C(S^1)^{(n)}$ has $\mathrm{prop} = 2$
(their Proposition~4.2), while the circulant matrix algebra
$C(C_m)$ has $\mathrm{prop} = 1$ trivially (it is its own
$C^*$-algebra).  Eqs.~\eqref{eq:prop_two}
(Proposition~\ref{prop:propagation_number}) and~\eqref{eq:prop_circulant}
(Corollary~\ref{cor:structural_specificity}) reproduce both axes
verbatim on $S^3$:
\[
\renewcommand{\arraystretch}{1.15}
\begin{array}{c|c|c}
 & \text{spectral truncation} & \text{circulant analog} \\ \hline
S^1 \text{ (Connes--vS)} & C(S^1)^{(n)},\ \mathrm{prop} = 2 &
   C(C_m),\ \mathrm{prop} = 1 \\
S^3 \text{ (this work)}  & \mathcal{O}_{n_{\max}},\ \mathrm{prop} = 2 &
   A_{\mathrm{circ}, N},\ \mathrm{prop} = 1 \text{ or } \infty \\
\end{array}
\]
The two-sided match foreclosing the natural objection that
$\mathrm{prop} = 2$ might be a generic feature of any finite
truncation of $S^3$.  Together with the SO(4)-isometry-preservation
property of Sec.~\ref{sec:axiom_audit} below (the property Connes
\& van~Suijlekom themselves identify as distinguishing spectral
truncation from circulant discretization), this is the strongest
quantitative alignment between GeoVac and the published
spectral-truncation framework currently known.
\end{remark}

\subsection{The Hilbert space $\mathcal{H}_{\mathrm{GV}}$}

Let $V_{\mathrm{Dirac}}$ denote the set of Dirac-graph nodes at the
matching truncation, labeled by Camporesi--Higuchi triples
$(n_{\mathrm{D}}, \kappa, m_j)$ with $0 \le n_{\mathrm{D}} \le
n_{\max} - 1$ (the standard convention $n_{\mathrm{Fock}} = n_{\mathrm{CH}}
+ 1$ from \texttt{geovac/dirac\_matrix\_elements.py}), Dirac--Coulomb
relativistic angular momentum $\kappa \in \{-l-1, +l\}$ for each
$0 \le l \le n_{\mathrm{D}}$, and total angular momentum projection
$m_j \in \{-j, \ldots, j\}$ where $j = |\kappa| - 1/2$.  The total
node count is
\[
   N_{\mathrm{Dirac}}(n_{\max}) = \sum_{n=0}^{n_{\max}-1}
     g_n^{\mathrm{D}} = \sum_{n=0}^{n_{\max}-1} 2(n+1)(n+2)
   = \tfrac{2}{3} n_{\max}(n_{\max}+1)(n_{\max}+2).
\]

\begin{definition}[Spinor space on the Dirac graph]
\label{def:H_GV}
$\mathcal{H}_{\mathrm{GV}}^{(n_{\max})} := \mathbb{C}^{V_{\mathrm{Dirac}}}$,
with the standard Hermitian inner product on $\mathbb{C}^{N_{\mathrm{Dirac}}}$.
\end{definition}

The faithful $*$-representation of $\mathcal{A}_{\mathrm{GV}}$ on
$\mathcal{H}_{\mathrm{GV}}$ is by the natural pullback: for
$f \in \mathcal{A}_{\mathrm{GV}}$ and $\psi \in \mathcal{H}_{\mathrm{GV}}$,
\begin{equation}
   (\pi(f) \psi)(n_{\mathrm{D}}, \kappa, m_j)
   := f(n_{\mathrm{Fock}}(n_{\mathrm{D}}), l(\kappa), m_l(\kappa, m_j))
   \cdot \psi(n_{\mathrm{D}}, \kappa, m_j),
   \label{eq:representation}
\end{equation}
where the maps $n_{\mathrm{Fock}}, l, m_l$ are the standard Dirac-to-Fock
quantum-number bridges:
$n_{\mathrm{Fock}}(n_{\mathrm{D}}) = n_{\mathrm{D}} + 1$,
$l(\kappa) = -1 - \kappa$ if $\kappa < 0$ and $l(\kappa) = \kappa$ if
$\kappa > 0$, and $m_l$ via the standard
$|j, m_j\rangle = \langle l, m_l; \tfrac{1}{2}, m_s | j, m_j \rangle
|l, m_l\rangle |s, m_s\rangle$ Clebsch--Gordan decomposition.
Each $\pi(f)$ acts diagonally; in the spinor basis it is the diagonal
matrix $\mathrm{diag}(f(v_{\mathrm{Fock}}(\,\cdot\,)))$.  The
representation is faithful: $\pi(f) = 0 \Leftrightarrow f \equiv 0$
(every Fock node has at least one Dirac node above it for $n_{\max}
\ge 1$).

\subsection{The Dirac operator $D_{\mathrm{GV}}$}

The continuum Camporesi--Higuchi Dirac operator on the unit
three-sphere has spectrum $|\lambda_n^{\mathrm{D}}| = n + 3/2$ with
degeneracy $g_n^{\mathrm{D}} = 2(n+1)(n+2)$~\cite{camporesi_higuchi1996}.
We promote this to a finite operator on $\mathcal{H}_{\mathrm{GV}}$
in two equivalent ways and verify their agreement in the truncation.

\begin{definition}[Spectral form of $D_{\mathrm{GV}}$]
\label{def:D_GV_spectral}
$D_{\mathrm{GV}}^{\mathrm{spec}}$ is the diagonal operator on
$\mathcal{H}_{\mathrm{GV}}$ with eigenvalues
$\lambda_n^{\mathrm{D},\pm} = \pm (n_{\mathrm{D}} + 3/2)$
on the eigenstate $|n_{\mathrm{D}}, \kappa, m_j\rangle$ with sign
$+$ for $\kappa < 0$ and $-$ for $\kappa > 0$ (the standard
``large-component'' / ``small-component'' sign choice of the
Camporesi--Higuchi convention~\cite{camporesi_higuchi1996}).
\end{definition}

\begin{definition}[Graph form of $D_{\mathrm{GV}}$]
\label{def:D_GV_graph}
$D_{\mathrm{GV}}^{\mathrm{graph}}$ is the discrete first-order
hopping operator on the Dirac graph defined by the Camporesi--Higuchi
nearest-neighbor edge set with weights chosen so that the spectrum
matches Definition~\ref{def:D_GV_spectral} in the truncation.  The
edge set and weights are documented in
\texttt{geovac/dirac\_matrix\_elements.py}; the explicit form, written
in the $(n_{\mathrm{D}}, \kappa, m_j)$ basis, is the
Lichnerowicz-derived first-order operator with $\kappa$-coupled
adjacent-shell amplitudes.
\end{definition}

\begin{proposition}[Equivalence of spectral and graph forms]
\label{prop:D_equiv}
For all $n_{\max} \ge 1$, $D_{\mathrm{GV}}^{\mathrm{spec}}
= U^{\dagger} D_{\mathrm{GV}}^{\mathrm{graph}} U$ for the unitary
$U$ that diagonalizes $D_{\mathrm{GV}}^{\mathrm{graph}}$.  The
agreement is exact in finite arithmetic at every truncation; in
particular the spectrum, multiplicities, and trace invariants
$\mathrm{Tr}\,(D_{\mathrm{GV}})^{2k}$ match the truncated
Camporesi--Higuchi continuum values.
\end{proposition}

\begin{proof}[Proof sketch.]
Both operators are constructed to have the truncated
Camporesi--Higuchi spectrum by definition.  The equivalence is the
statement that there exists an orthonormal basis of
$\mathcal{H}_{\mathrm{GV}}$ in which both are diagonal, which holds
because $D_{\mathrm{GV}}^{\mathrm{graph}}$ is self-adjoint by
construction (Hermitian conjugate edge weights) and has the prescribed
spectrum.  Numerical agreement at $n_{\max} \le 6$ is verified by
existing tests in \texttt{tests/test\_dirac\_matrix\_elements.py} (108
tests) and \texttt{tests/test\_qed\_self\_energy.py}.
\end{proof}

We use the spectral form $D_{\mathrm{GV}}^{\mathrm{spec}}$ (henceforth
$D_{\mathrm{GV}}$) for axiom verification because it is closed-form;
the graph form is what one uses for finite-truncation computations of
the spectral action.  The $\pi$-free certificate of the Dirac sector
(Paper~28~\S graph\_native\_qed.6) applies to both: every eigenvalue
is a half-integer rational, every degeneracy is a positive integer.

\subsection{Bounded commutators}

\begin{proposition}[Bounded commutators]
\label{prop:bounded_commutators}
For every $a \in \mathcal{A}_{\mathrm{GV}}$, the commutator
$[D_{\mathrm{GV}}, \pi(a)]$ is a bounded operator on
$\mathcal{H}_{\mathrm{GV}}$.
\end{proposition}

\begin{proof}
$D_{\mathrm{GV}}$ has finite-dimensional domain
$\mathcal{H}_{\mathrm{GV}}$ at every $n_{\max}$, hence bounded.
$\pi(a)$ is a diagonal multiplication operator, also bounded.  Any
commutator of bounded operators is bounded.
\end{proof}

This is the trivial but necessary check: in finite truncation, every
operator is bounded.  The non-trivial content---that
$\|[D_{\mathrm{GV}}, \pi(a)]\|$ is controlled uniformly in
$n_{\max}$---is what makes the continuum limit a bona-fide spectral
triple in the sense of Connes; for a Lipschitz function $a$, the
commutator norm equals the Lipschitz constant on the graph metric, in
exact analogy with the smooth case.

\subsection{The triple, summarized}

Putting Definitions~\ref{def:A_GV},~\ref{def:H_GV},~\ref{def:D_GV_spectral}
and Proposition~\ref{prop:bounded_commutators} together:

\begin{theorem}[The GeoVac spectral triple]
\label{thm:GV_triple}
For every $n_{\max} \ge 1$, the data
$(\mathcal{A}_{\mathrm{GV}}^{(n_{\max})},
\mathcal{H}_{\mathrm{GV}}^{(n_{\max})},
D_{\mathrm{GV}}^{(n_{\max})})$ is a finite spectral triple in the
sense of~\cite{connes1995}: $\mathcal{A}_{\mathrm{GV}}$ is unital and
involutive, $\pi: \mathcal{A}_{\mathrm{GV}} \to
\mathcal{B}(\mathcal{H}_{\mathrm{GV}})$ is a faithful
$*$-representation, $D_{\mathrm{GV}}$ is self-adjoint with bounded
commutators.  The continuum limit $n_{\max} \to \infty$ recovers the
canonical spectral triple of the unit Riemannian spin three-sphere
$S^3$ with the Camporesi--Higuchi Dirac operator.
\end{theorem}

The continuum-limit clause is the spectral-truncation statement of
Connes--van~Suijlekom~\cite{connes_vs2021,hekkelman2022,hekkelman2024}:
the GeoVac triples form a Cauchy sequence in the truncation order whose
limit is the standard $S^3$ spectral triple.

% ===================================================================
\section{Axiom Audit}
\label{sec:axiom_audit}
% ===================================================================

We check each Connes axiom in turn, distinguishing what holds
rigorously at finite $n_{\max}$ from what holds only in the continuum
limit.

\subsection{Order zero (involutive faithful representation)}

Holds rigorously by Definitions~\ref{def:A_GV}--\ref{def:H_GV} and
Eq.~\eqref{eq:representation}.  Verified in
\texttt{tests/test\_dirac\_matrix\_elements.py}.

\subsection{Bounded commutators (regularity)}

Holds rigorously at every $n_{\max}$ by
Proposition~\ref{prop:bounded_commutators}.

\subsection{Compact resolvent / finite summability}

In the finite-$n_{\max}$ truncation $\mathcal{H}_{\mathrm{GV}}$ is
finite-dimensional so this is automatic.  The continuum statement is
that $D_{\mathrm{GV}}$ has compact resolvent and is $3^+$-summable
(Dixmier-trace summable in dimension $3+\epsilon$); this is the
Camporesi--Higuchi summability of the continuum operator on $S^3$ and
is a standard result~\cite{camporesi_higuchi1996,connes1995}.

\subsection{Order one}

\begin{proposition}[Order one]
\label{prop:order_one}
For every $a, b \in \mathcal{A}_{\mathrm{GV}}$,
$[[D_{\mathrm{GV}}, \pi(a)], \pi(b)^{\circ}] = 0$ holds
trivially because the right action $\pi(b)^{\circ}$ commutes with the
left action $\pi(a)$ (the algebra is commutative) and with the
diagonal-in-spectral-basis form of $D_{\mathrm{GV}}$.
\end{proposition}

This is the standard observation that the order-one condition is
automatic for any commutative spectral triple.  Order-one becomes
non-trivial only when one extends to an almost-commutative triple
$\mathcal{A}_{\mathrm{GV}} \otimes M_n(\mathbb{C})$ with non-abelian
fiber, at which point it constrains the form of fluctuations of $D$.
This is exactly the setting where Papers~25 ($n=1$, $U(1)$ fiber) and
30 ($n=2$, $\mathrm{SU}(2)$ fiber) live; we treat them in
Sec.~\ref{sec:subsectors}.

\subsection{Real structure}

$S^3$ is odd-dimensional; in dimension~$n$ the canonical
$\mathrm{KO}$-dimension is~$n$ mod~$8$ for a Riemannian spin manifold,
and the sign table~\cite{connes1995,vansuijlekom_book2015} for $n=3$
gives
\[
   \epsilon = -1, \quad \epsilon' = +1, \quad \epsilon'' = ~\text{(none)}~
   ~~\text{(no $\gamma$ in odd dimension).}
\]
The sign $\epsilon = -1$ reflects the fact that the 2-dimensional
spinor representation of $\mathrm{Spin}(3) = \mathrm{SU}(2)$ is
quaternionic (pseudoreal): the standard charge-conjugation operator
on this representation squares to $-\mathbb{1}$, equivalently
$\mathbb{H} \otimes_{\mathbb{R}} \mathbb{C} \cong M_2(\mathbb{C})$
acts on a complex 2-dimensional space with antiunitary $J$ satisfying
$J^2 = -\mathbb{1}$.  This is the same sign that appears in the
canonical Lorentzian Dirac formalism in $3+1$ dimensions for the
charge-conjugation matrix~$C$.

\begin{proposition}[Real structure on $\mathcal{H}_{\mathrm{GV}}$]
\label{prop:reality}
There exists an antilinear isometry $J_{\mathrm{GV}}:
\mathcal{H}_{\mathrm{GV}} \to \mathcal{H}_{\mathrm{GV}}$ with
$J_{\mathrm{GV}}^2 = -\mathbb{1}$ and $J_{\mathrm{GV}} D_{\mathrm{GV}}
= +D_{\mathrm{GV}} J_{\mathrm{GV}}$, namely the Camporesi--Higuchi
charge-conjugation operator restricted to the truncation.
\end{proposition}

\begin{proof}[Construction.]
The continuum charge conjugation $J_{S^3}$ on the
Camporesi--Higuchi spinor bundle satisfies the listed signs by
construction~\cite{camporesi_higuchi1996,connes1995}.  Restriction to
the truncation $\mathcal{H}_{\mathrm{GV}}^{(n_{\max})}$ commutes with
the spectral truncation because $D_{\mathrm{GV}}$ acts diagonally in
the spectral basis and $J_{S^3}$ permutes the spectral basis (sending
$|n_{\mathrm{D}}, \kappa, m_j\rangle \mapsto
|n_{\mathrm{D}}, -\kappa, -m_j\rangle$ up to a phase, whose
antilinear-squaring gives $-\mathbb{1}$ on the quaternionic 2-spinor
representation).  The restriction therefore inherits the signs.
\end{proof}

\subsection{$\mathbb{Z}_2$-grading}

\begin{proposition}[Odd triple]
\label{prop:odd}
The GeoVac spectral triple is \emph{odd}: there is no chirality
operator $\gamma$ on $\mathcal{H}_{\mathrm{GV}}$ commuting with
$\pi(\mathcal{A}_{\mathrm{GV}})$ and anticommuting with
$D_{\mathrm{GV}}$.
\end{proposition}

\begin{proof}
$S^3$ is odd-dimensional and the Camporesi--Higuchi Dirac operator on
$S^3$ has no chirality operator in the standard sense.  The truncation
inherits the same statement.
\end{proof}

This is correct, not a defect.  Odd-dimensional spectral triples do
not carry a $\mathbb{Z}_2$-grading; their spectral action is
constructed from the absolute value $|D|$ rather than from $D^2$
projected onto chirality eigenspaces, and this is precisely the
structure that Paper~28's QED-on-$S^3$ computations use.

\subsection{Summary}

\begin{table}[h]
\centering
\begin{tabular}{lll}
\toprule
Axiom & Status at finite $n_{\max}$ & Status in continuum limit \\
\midrule
Order zero (faithful $*$-rep) & Holds rigorously & Holds rigorously \\
Bounded commutators           & Holds rigorously & Holds rigorously \\
Compact resolvent             & Auto (finite dim) & Holds (CH summability) \\
Finite summability            & Auto              & $3^+$-summable \\
Order one                     & Holds trivially   & Holds trivially \\
Real structure $J$            & Holds (CH conjugation) & Holds \\
$J^2 = -\mathbb{1}$           & Holds (quaternionic) & Holds \\
$JD = +DJ$                    & Holds             & Holds \\
$\mathbb{Z}_2$-grading        & Absent (odd)      & Absent (odd) \\
\bottomrule
\end{tabular}
\caption{Axiom audit for the GeoVac triple
$(\mathcal{A}_{\mathrm{GV}}, \mathcal{H}_{\mathrm{GV}}, D_{\mathrm{GV}})$.
All standard axioms are satisfied; the triple is real, odd, and of
$\mathrm{KO}$-dimension~$3$ in the continuum limit.}
\label{tab:axiom_audit}
\end{table}

\begin{theorem}[Rigorous identification]
\label{thm:axiom_complete}
The continuum limit of the GeoVac spectral triple is the canonical
spectral triple of the unit Riemannian spin three-sphere with the
Camporesi--Higuchi Dirac operator.  In particular it is real, odd, and
of $\mathrm{KO}$-dimension~$3$.
\end{theorem}

This is a structural theorem; no fitting parameters or empirical
inputs.  The audit in Table~\ref{tab:axiom_audit} establishes that the
spectral-triple framing is not metaphorical.

% ===================================================================
\section{Sub-sector Identification: Papers 25, 28, 30, 31}
\label{sec:subsectors}
% ===================================================================

We now read the major framework papers as projections of the GeoVac
triple.

\subsection{Paper 28 (QED on $S^3$): the spectral action}

Paper~28~\cite{paper28} computes Tr~$f(D^2/\Lambda^2)$ and related
spectral sums on the Dirac operator $D_{\mathrm{GV}}$ in the continuum
limit, with $\mathrm{SO}(4)$ vertex selection rules from the
Hopf-bundle decomposition.  In the language of Sec.~\ref{sec:connes_data}:
\begin{itemize}\setlength{\itemsep}{0pt}
\item Theorem~T9 (Paper~28 Sec.~T9): the spectral zeta of $D^2$ is a
two-term polynomial in $\pi^2$ with rational coefficients at every
integer $s$.  In spectral-triple language: the heat-kernel
coefficients $a_k(D^2)$ on $S^3$ saturate at $k=2$ (Paper~28's
``SD two-term exactness theorem'').
\item Self-energy structural zero (Paper~28 Theorem~4):
$\Sigma(n_{\mathrm{ext}}=0) = 0$ on the continuum spectral sum,
identically by vertex parity.  This is the standard heat-kernel
bosonic sector cancellation specialized to the Camporesi--Higuchi
spectrum.
\item Vector-photon QED on $S^3$ (Paper~28 \S vector\_photon\_qed):
adding photon $(q, m_q)$ labels with Wigner 3$j$ vertex coupling
recovers 7/8 selection rules for scalar electrons and 8/8 for Dirac
electrons.  The 1/(4$\pi$) per-loop calibration is the $S^2$ Weyl
exchange constant of the Hopf base.
\end{itemize}

In the spectral-triple language Paper~28 is the spectral-action
computation $\mathrm{Tr}\,f(D_{\mathrm{GV}}^2/\Lambda^2)$ at one loop;
the Camporesi--Higuchi spectrum is the input data of $D_{\mathrm{GV}}$;
the SO(4) selection rules are the harmonic decomposition of the spinor
bundle.  Paper~28 is not a free-standing observation about $S^3$; it is
the spectral action of the present triple.

\subsection{Paper 25 (Hopf graph as lattice gauge): the Cartan torus}

Paper~25~\cite{paper25} introduces the signed node-edge incidence
matrix $B$ of the Fock graph and identifies $L_0 = BB^{\top}$ as the
matter propagator and $L_1 = B^{\top}B$ as the discrete
Hodge-1 Laplacian.  In the spectral-triple language: extend
$\mathcal{A}_{\mathrm{GV}}$ to
$\mathcal{A}_{\mathrm{GV}} \otimes \mathbb{C}$ (trivial $M_1$ fiber)
and consider inner fluctuations of $D_{\mathrm{GV}}$ in the abelian
$U(1)$ inner-derivation algebra.  The fluctuated Dirac operator
$D_{\mathrm{GV}} + A$, with $A \in \Omega_D^1$ a self-adjoint 1-form
in the $D_{\mathrm{GV}}$-calculus, has $L_1$ as its quadratic kinetic
term.  This is the Cartan-subalgebra (maximal-torus) sector.

Paper~25's $L_1$ is the $U(1)$-fluctuated kinetic term of the GeoVac
spectral action; its assignment as ``photon propagator'' in Paper~25
is exactly the spectral-action reading.

\subsection{Paper 30 (SU(2) Wilson gauge): the full non-abelian sector}

Paper~30~\cite{paper30} builds the $\mathrm{SU}(2)$ Wilson lattice
gauge theory on the same graph and proves: (a) the maximal-torus
reduction reproduces Paper~25's $U(1)$ action exactly; (b) the
weak-coupling kinetic term equals $L_1$ acting on
$\mathfrak{su}(2)$-valued 1-cochains.  In the spectral-triple
language: extend $\mathcal{A}_{\mathrm{GV}}$ to
$\mathcal{A}_{\mathrm{GV}} \otimes M_2(\mathbb{C})$ (a non-abelian
$M_2$ fiber, equivalently a rank-1 non-abelian almost-commutative
extension of the GeoVac triple).  The order-one condition for this
extension is non-trivial and constrains the form of fluctuations to be
$\mathfrak{su}(2)$-valued connections, exactly Paper~30's link
variables.  The Wilson plaquette action is the lattice-gauge
implementation of the spectral-action curvature term in this
extension.

Paper~30 is the rank-1 non-abelian almost-commutative extension of the
GeoVac triple; Paper~25 is the abelian projection of Paper~30;
together they exhibit the standard Cartan torus / full group structure
of the gauge sector of the present triple.

\subsection{Paper 31 (universal/Coulomb partition): the
$\mathcal{A}/D$ split}

Paper~31~\cite{paper31} partitions GeoVac results into a universal
sector that depends only on the algebra $\mathcal{A}$ (angular
sparsity, composed scaling, selection rules) and a potential-specific
sector that depends on the Dirac operator $D$ (spectrum, calibration
$\kappa$, transcendental content, rigidity, $\alpha$ conjecture).  In
the spectral-triple language this is exactly the $\mathcal{A}/D$ split:
properties of $\mathcal{A}_{\mathrm{GV}}$ alone (rank, dimension,
representation content) transfer to any spectral triple with the same
algebra; properties of $D_{\mathrm{GV}}$ specifically depend on the
choice of Dirac operator.

Paper~31 is the structural reading of Paper~22 (angular sparsity
theorem, an $\mathcal{A}$-statement) and Paper~23 (nuclear shell
model, a different $D$-choice on the same $\mathcal{A}$): when one
keeps $\mathcal{A}$ fixed and varies $D$, the universal results
transfer; when one varies $\mathcal{A}$ (say, replace $S^3$ with
$S^5$, Paper~24), the entire structural picture reorganizes.

\subsection{Summary of sub-sector identification}

\begin{table}[h]
\centering
\begin{tabular}{lll}
\toprule
Paper & Object in this paper's language & Sector \\
\midrule
Paper~7 & $\mathcal{A}_{\mathrm{GV}}$ continuum limit & Manifold sector \\
Paper~22 & Selection rules from $\mathcal{A}_{\mathrm{GV}}$ & $\mathcal{A}$ sector (universal) \\
Paper~23 & $D$-replacement: nuclear shell model & Different $D$, same $\mathcal{A}$ \\
Paper~24 & Different $\mathcal{A}$: $S^5$ holomorphic & Distinct triple (HO) \\
Paper~25 & $U(1)$-fluctuated $D_{\mathrm{GV}} + A$ & Cartan-torus inner fluctuation \\
Paper~28 & Tr~$f(D_{\mathrm{GV}}^2/\Lambda^2)$ & Spectral action (one loop) \\
Paper~30 & $\mathrm{SU}(2)$-fluctuated, $M_2$ fiber & Non-abelian extension \\
Paper~31 & $\mathcal{A}/D$ split & Structural partition \\
\bottomrule
\end{tabular}
\caption{Sub-sector identification of the major framework papers as
projections of the GeoVac spectral triple.}
\label{tab:subsectors}
\end{table}

% ===================================================================
\section{The $K = \pi(B + F - \Delta)$ Formula in Spectral-Action
Language}
\label{sec:K_decomposition}
% ===================================================================

Paper~2~\cite{paper2} observed the numerical coincidence
$K = \pi(B + F - \Delta) = 1/\alpha$ at $8.8 \times 10^{-8}$, with
\[
   B = 42, \quad F = \tfrac{\pi^2}{6}, \quad \Delta = \tfrac{1}{40},
\]
and the Phase~4B--4I sprint series identified three independent
spectral homes for the three pieces (CLAUDE.md~\S 2):
\begin{itemize}\setlength{\itemsep}{0pt}
\item $B$ is the finite Casimir trace at $m=3$:
$B = \sum_{n=1}^{3} (2l+1)l(l+1)|_{n,l} = 42$ (Paper~2 Eq.~17).
\item $F = D_{n^2}(d_{\max} = 4) = \zeta_R(2)$, the Fock-degeneracy
Dirichlet series at the packing exponent (Phase~4F $\alpha$-J).
\item $\Delta^{-1} = g_3^{\mathrm{D}} = 40$, the third single-chirality
Camporesi--Higuchi Dirac mode degeneracy (Phase~4H SM-D).
\end{itemize}

In spectral-triple language, $B$, $F$, $\Delta$ live in three different
sectors of the same triple:

\begin{itemize}\setlength{\itemsep}{2pt}
\item $B$ is a \emph{finite scalar Casimir trace} on the truncated
algebra: in the language of Sec.~\ref{sec:axiom_audit}, it is
$\mathrm{tr}_{\mathcal{A}_{\mathrm{GV}}}(C \cdot \mathbb{1})$ for $C$
the second Casimir of $\mathrm{SO}(4)$ at $n_{\max} = 3$.
\item $F$ is an \emph{infinite scalar spectral zeta}: it equals
$\zeta_{\Delta_{S^3}}(2)$ in the continuum limit, where
$\Delta_{S^3}$ is the scalar Laplace--Beltrami operator on $S^3$
(equivalently, the eigenvalue-weighted Dirichlet of the Fock
degeneracy).
\item $\Delta^{-1}$ is a \emph{single-level Dirac degeneracy}: it
equals $g_3^{\mathrm{D}}$, the multiplicity of the third
Camporesi--Higuchi single-chirality eigenvalue
$|\lambda_3^{\mathrm{D}}| = 9/2$.
\end{itemize}

These three are categorically different objects: a finite
combinatorial trace, an infinite analytic Dirichlet series, and a
single-level multiplicity.  They sit in different sectors of the
spectral triple (scalar/spinor; finite/infinite; trace/zeta/dimension)
and Phases 4B--4I exhausted nine attempts to find a common spectral
generator that produces all three simultaneously, without success.

\begin{observation}[Three-sector composition]
\label{obs:three_sector}
The Paper~2 combination rule $K = \pi(B + F - \Delta)$ is a sum across
three structurally distinct sectors of the GeoVac spectral triple.
Each piece has been individually identified as a clean spectral
invariant; the combination rule itself remains a numerical observation.
\end{observation}

This rephrasing is what working hypothesis WH5 in the project register
is built on: $\alpha$ is the conversion factor between three projection
regimes that do not share a generator.  The spectral-triple framing
does not derive $\alpha$; it makes precise where the obstruction is.

\begin{remark}[Sprint A, $\pi^3 \alpha^3$ residual]
\label{rem:sprint_A}
Sprint~A (April~2026, CLAUDE.md~\S 1.7) verified at 80~dps that the
post-cubic residual $K - 1/\alpha - \alpha^2 = 1.2079 \times 10^{-5}$
matches $\pi^3 \alpha^3 = 1.2049 \times 10^{-5}$ to within 0.25\%.
The structural reading is $\pi^3 = \mathrm{Vol}(S^5)$, hinting at a
contribution from the Bargmann--Segal $S^5$ HO sector
(Paper~24~\cite{paper24}) at higher order.  This is a structural
hint, not a derivation; the cross-check (Paper~24's HO triple is a
distinct spectral triple, not a sub-sector of the GeoVac triple)
shows that the $\pi^3 \alpha^3$ identification, if real, is a
\emph{cross-triple} relation and not internal to the GeoVac
spectral-action expansion.
\end{remark}

% ===================================================================
\section{The Coulomb/HO Asymmetry in Spectral-Triple Language}
\label{sec:coulomb_ho}
% ===================================================================

Paper~31~\cite{paper31} formalized a three-layer asymmetry between the
Coulomb spectral triple of the present paper and the harmonic
oscillator spectral triple of Paper~24~\cite{paper24}, summarized
here:

\begin{table}[h]
\centering
\small
\begin{tabular}{lll}
\toprule
Layer & Coulomb (this paper) & Harmonic oscillator (Paper~24) \\
\midrule
Manifold algebra $\mathcal{A}$ &
  Functions on $S^3$ Fock graph &
  Functions on $S^5$ holomorphic Hardy graph \\
Operator order &
  Second-order (Laplace--Beltrami) &
  First-order (complex Euler) \\
Spectrum &
  Quadratic $|\lambda_n| = n^2 - 1$ &
  Linear $E_N = N + 3/2$ \\
Calibration constant &
  $\pi$ (Hopf measure $\mathrm{Vol}(S^2)/4$) &
  None ($\pi$-free in exact arithmetic) \\
Wilson gauge content &
  Full $\mathrm{SU}(2)$ (Paper~30) &
  $U(1)$ only (Sprint~5 negative) \\
$L_0$ matter propagator role &
  Yes (eigenvalues = spectrum) &
  No (only adjacency, not spectrum) \\
\bottomrule
\end{tabular}
\caption{The Coulomb/HO asymmetry from Paper~31, in spectral-triple
language.  The two systems are distinct spectral triples sharing only
the algebra type (functions on a graph) and the trace structure of the
spectral action.}
\label{tab:coulomb_ho}
\end{table}

In spectral-triple language the asymmetry is structural: the Coulomb
case is a Riemannian-spin spectral triple of $\mathrm{KO}$-dimension~3,
the HO case is a spectral triple of complex Euler type on the
holomorphic Hardy space of $S^5$.  These are different objects even
though both are ``almost-commutative'' in the broad sense.  Calibration
$\pi$ enters in the Coulomb case through the Hopf-bundle measure
$\mathrm{Vol}(S^2) / 4$ in the spectral-action coefficients; it does
not enter the HO case because the HO spectral action is built from a
linear (first-order) operator whose heat kernel does not produce $\pi$
at the rational level of the truncation~\cite{paper24}.

\begin{observation}[Coulomb-specificity of $\pi$]
\label{obs:pi_coulomb}
The presence of $\pi$ in the spectral-action coefficients of QED is
specific to the Coulomb-projected $S^3$ triple of this paper, not a
universal feature of GeoVac-style discretizations.  The HO triple is
$\pi$-free at the discrete level.  This is the structural answer to
``why does QED need $\pi$''---it needs $\pi$ because it lives on the
Coulomb-projection $S^3$ triple, where the calibration is the
$S^2$ Hopf-base Weyl constant.
\end{observation}

% ===================================================================
\section{The $\pi$-Source Case-Exhaustion Theorem}
\label{sec:pi_source_theorem}
% ===================================================================

The empirical principle of Paper~35~\cite{paper35} (Prediction~1,
graduated to load-bearing status on a cumulative 208/208 confirmation
across sprints KG and TX-B, May~2026) states that a GeoVac observable
contains $\pi$ if and only if its evaluation includes a continuous
integration over a temporal or spectral parameter promoted from the
discrete graph spectrum.  Read in the spectral-triple language of
\S\ref{sec:construction}--\S\ref{sec:axiom_audit}, this principle has
a structural underpinning: the only places in the GeoVac framework
where such continuous integrations occur are anchored to three named
mechanisms of the present triple.  We state and prove the
case-exhaustion.

\begin{theorem}[$\pi$-Source Case-Exhaustion]
\label{thm:pi_source_case_exhaustion}
Let $\mathcal{C} = (P_{i_k} \circ \cdots \circ P_{i_1})$ be any finite
composition of projections drawn from the Paper~34 list~\cite{paper34}
\S\,III.1--\S\,III.15, applied to a Layer-1 datum on
$(\mathcal{A}_{\mathrm{GV}}, \mathcal{H}_{\mathrm{GV}},
D_{\mathrm{GV}})$ or its Bargmann--Segal sibling triple.  The
following are equivalent:
\begin{enumerate}
\setlength\itemsep{0pt}
\item[(i)] The converged value $\mathcal{C}(\mathrm{datum})$ contains a
transcendental of the form $\pi$, $\sqrt{\pi}$, $\pi^k$ for $k \in
\mathbb{Q}$, or $\pi$ in a denominator.
\item[(ii)] The chain $\mathcal{C}$ contains at least one projection
$P_{i_j}$ whose evaluation pipeline includes a continuous integration
over a parameter promoted from the discrete graph spectrum
(heat-kernel time, Matsubara frequency, Mellin contour, Hopf $S^2$-base
angular measure, $\mathrm{SU}(N)$ Haar measure, or
analytic-continuation parameter).
\item[(iii)] The chain $\mathcal{C}$ engages at least one of three
spectral-triple mechanisms:
\begin{itemize}\setlength\itemsep{0pt}
\item[\textbf{M1}] (Hopf-base measure):\ introduces $\pi =
\mathrm{Vol}(S^2)/4$ via the Hopf bundle $S^3 \to S^2 \times S^1$,
or its higher-rank analogues~\cite{paper25,paper30}.
\item[\textbf{M2}] (Seeley--DeWitt heat-kernel coefficient):\
introduces $\sqrt{\pi}$ via heat-kernel coefficients on the
Camporesi--Higuchi Dirac spectrum on unit~$S^3$ (Paper~28
Theorem~T9~\cite{paper28}: $\zeta_{D^2}(s)$ is a two-term polynomial
in $\pi^2$ with rational coefficients at every integer $s \ge 1$).
The sub-mechanism for the spinor sector is the half-integer Hurwitz
$\zeta(s, 3/2)$.
\item[\textbf{M3}] (vertex-topology Dirichlet $L$-content):\
introduces Catalan~$G = \beta(2)$ and Dirichlet $\beta(s)$ via
quarter-integer Hurwitz shifts $\zeta(s, 3/4)$, $\zeta(s, 5/4)$ at
two-loop sunset vertices (Paper~28 Theorem~3:\
$D_{\mathrm{even}}(s) - D_{\mathrm{odd}}(s) = 2^{s-1}(\beta(s) -
\beta(s-2))$).
\end{itemize}
\end{enumerate}
\end{theorem}

\begin{proof}[Proof sketch]
(i)~$\Leftrightarrow$~(ii) is Paper~35 Prediction~1, verified across
208 individual checks (sprints KG, TX-B, May 2026~\cite{paper35}).

(ii)~$\Rightarrow$~(iii):\ each of the fifteen Paper~34 projections is
either $\pi$-free at the projection step (projections~1,~3,~4,~5,~8,
9,~12,~13,~14: Fock conformal, Bargmann--Segal, Stereographic, Sturmian,
Wigner~$3j$, Wigner~$D$, Molecule-frame, Drake--Swainson, Rest-mass) or
$\pi$-bearing through a continuous integration that reduces to one of
M1,~M2,~M3 by direct identification (Hopf~bundle, Spectral~action,
Camporesi--Higuchi spinor lift, Wilson plaquette, Vector-photon promotion,
Observation/temporal-window).  The seven $\pi$-free projections live in
either Paper~31's universal sector (algebra-only, transfers to any
spherical-fermion system) or Paper~22's potential-independent angular
sparsity tier; neither sector contains $\pi$.  The eight $\pi$-bearing
projections are tagged in
\texttt{debug/track\_ts\_e1\_theorem\_promotion\_memo.md}~\S2 to one or
more of M1/M2/M3, with closed-form identification for each (Paper~25
for M1, Paper~28 T9 theorem for M2, Paper~28 Theorem~3 for M3).

(iii)~$\Rightarrow$~(i) is closed-form for each mechanism:\ M1's
$\mathrm{Vol}(S^2)/4 = \pi$ identity, M2's $a_0(S^3) = \sqrt{\pi}$
heat-kernel identity, M3's $\beta(s) - \beta(s-2)$ closed form via
the $\chi_{-4}$ Dirichlet character.  Each forces $\pi$ to appear in
any chain that engages it.
\end{proof}

\begin{remark}[Joint engagement]
M1, M2, M3 are not set-theoretically disjoint as constructions:\ the
spectral action invokes both M1 (Hopf-base measure on the spinor
bundle) and M2 (Seeley--DeWitt coefficients) simultaneously, and the
vacuum-polarization coefficient $\Pi = 1/(48\pi^2)$ (Paper~28~\S\,V)
contains one factor of $\pi$ from each mechanism.  The theorem requires
only that \emph{at least one} of the three be engaged; multiple
engagement is the typical case at higher loop order.  The mechanisms
are role-disjoint in the Paper~18~\S\,III.7 Mellin-engine sense
(M1~$=$ packing-tier $\omega_d$, M2~$=$ heat-kernel via Mellin/Gaussian,
M3~$=$ vertex-parity quarter-integer Hurwitz character mechanism).
The Track~TS-E3 sprint (May~2026) further sharpened this by showing
that M1, M2, M3 are all sub-cases of a single master Mellin-engine
mechanism $\mathrm{Tr}(D^k \cdot e^{-tD^2})$ on a compact Riemannian
manifold, with $k \in \{0, 1, 2\}$ selecting the sub-mechanism.
\end{remark}

\begin{remark}[Status of $K = \pi(B + F - \Delta)$]
\label{rem:K_under_theorem}
Paper~2's combination $K = \pi(B+F-\Delta) \approx 1/\alpha$ matches
its target at $8.8 \times 10^{-8}$.  Under
Theorem~\ref{thm:pi_source_case_exhaustion}, $K$'s chain
(Fock~$\circ$~Hopf$\circ$~spinor) engages M1 (the explicit prefactor $\pi$,
identified by Sprint~A as $\mathrm{Vol}(S^2)/4$~\cite{paper2}) and M2
(in $F = \pi^2/6 = D_{n^2}(d_{\max})$ via the Mellin transform of
the scalar Fock-degeneracy).  The $\pi$-content is therefore consistent
with case-exhaustion.  The remaining anomaly is at the
\emph{depth-prediction} layer of Paper~34~\S\,VIII.3 (depth predicts
$\sim 3\%$ residual at three projections; $K$ matches at
$10^{-8}$): a quantitative-error puzzle, distinct from the
qualitative case-exhaustion of the present theorem, and
unaffected by it.
\end{remark}

\begin{remark}[Falsification target]
A counter-example to Theorem~\ref{thm:pi_source_case_exhaustion}
would be a sixteenth projection (or a finite chain of existing
projections) introducing $\pi$ via a fourth mechanism not in
\{M1, M2, M3\}.  Three natural candidates were investigated in
Track~TS-E3 (May~2026) and each REDUCED to combinations of the
three existing mechanisms:\ continuum-gravity coupling (M1+M2),
anomaly-class projections (M1+M2+M3 via the eta-invariant's
first-order Mellin), and non-trivial principal-bundle projections
(M1 normalisations against integer integrands).  The strongest
remaining computational falsification target is the discrete
first-Chern-class invariant on the Hopf $S^3$ graph at $n_{\max}=3$,
predicted integer-valued by the case-exhaustion; if a $\pi$ appears
unexpectedly, a fourth mechanism M4 would need to be named.
\end{remark}

\begin{remark}[Status of the Gromov--Hausdorff convergence roadmap]
\label{rem:gh_convergence_status}
The case-exhaustion theorem above pins the $\pi$-injection mechanisms
of the truncated GeoVac spectral triple at every finite cutoff
$n_{\max}$.  A complementary question is whether the truncated triple
\emph{converges} to the round-$S^3$ spectral triple in the
Gromov--Hausdorff sense as $n_{\max} \to \infty$, in the form
established for $\mathbb{T}^d$ by Leimbach--van~Suijlekom 2024 (Adv.
Math.\ 439, 109496).  Track~A of the WH1 sprint
(\texttt{debug/track\_ts\_a\_gh\_convergence\_memo.md}) lays out a
five-lemma proof shape for the SU(2)~$=S^3$ case via Peter--Weyl
analysis.  Two of the five lemmas are now closed: L1' (offdiag
Camporesi--Higuchi operator system substrate, every cross-pair
finite, sprint~R3.5, 2026-05-04;
\texttt{debug/wh1\_r35\_full\_dirac\_memo.md}); and L2 (existence of
a positive central spectral Fej{\'e}r kernel
$K_{n_{\max}} = Z_{n_{\max}}^{-1} |\sum_{j \le j_{\max}} \sqrt{2j+1}\,
\chi_j|^2$ on SU(2) with $\gamma_{n_{\max}} \to 0$, sprint~R2.5/L2,
2026-05-04; \texttt{debug/r25\_l2\_proof\_memo.md}).  The remaining
three lemmas (L3 Lipschitz bound, L4 Berezin reconstruction, L5
Latr{\'e}moli{\`e}re-propinquity assembly) are open at the date of
writing.  No paper edit beyond this remark; the GH-convergence
proof, when complete, is intended for a future companion paper
rather than insertion here.
\end{remark}

% -------------------------------------------------------------------
\subsection{The unified $U(1) \times SU(2) \times SU(3)$ gauge
content and its limits}
\label{sec:unified_gauge}
% -------------------------------------------------------------------

The case-exhaustion theorem of Sec.~\ref{sec:pi_source_theorem}
identifies M1 (Hopf-base measure) as the mechanism by which the
Wilson plaquette-Haar projection injects $\pi$ into spectral-triple
observables.  This subsection records a complementary structural
fact about the family of GeoVac sub-triples that engage M1: they
realise all three Standard-Model gauge groups as Wilson lattice
constructions, on three structurally distinct sub-manifolds, with a
single unified weak-coupling coefficient $1/(4\,N_c)$.  The
construction is honest about what it is and is not: a unified
\emph{catalogue} of three Wilson-Yang--Mills instances of the
Marcolli--van~Suijlekom gauge-network framework, not a unified
spectral triple in the Connes--Chamseddine
sense~\cite{chamseddine_connes2010}.

\subsubsection*{Three Wilson constructions, one kinetic coefficient}

Three independent sprints have produced gauge-invariant, finite,
Haar-normalisable non-abelian Wilson lattice gauge theories on
GeoVac sub-graphs:

\begin{itemize}\setlength{\itemsep}{2pt}
\item \emph{Paper~25}~\cite{paper25} constructs the abelian
$U(1)$ Wilson--Hodge structure on the Fock-projected $S^3$ Coulomb
graph (Paper~7), with phases $\chi_v$ on nodes and $A_e$ on edges
covariant under the standard incidence operator.
\item \emph{Paper~30}~\cite{paper30} constructs the non-abelian
$\mathrm{SU}(2)$ Wilson lattice gauge theory on the same
$S^3 = \mathrm{SU}(2)$ graph, with link variables $U_e \in
\mathrm{SU}(2)$, plaquettes from primitive non-backtracking walks,
and Haar measure $dU$.  Paper~30~\S 5.2 verifies the weak-coupling
kinetic coefficient symbolically and numerically.
\item \emph{Sprint~ST-SU3} (May~2026, modules
\texttt{geovac/su3\_wilson\_s5.py},
\texttt{tests/test\_su3\_wilson\_s5.py}, $39/39$ pass) constructs
the non-abelian $\mathrm{SU}(3)$ Wilson lattice gauge theory on the
Bargmann--Segal $S^5$ graph of Paper~24~\cite{paper24}, with link
variables $U_e \in \mathrm{SU}(3)$ in the fundamental representation
(independent of the shell index $N$), plaquettes from primitive
non-backtracking walks, and SU(3) Haar measure.  Sprint~ST-SU3~\S 4
verifies Theorem~2 (kinetic coefficient $1/12$) symbolically and
numerically at $N_{\max} = 2, 3$ ($7\,765$ plaquettes), with maximal
Cartan reduction to a coupled $U(1) \times U(1)$ theory verified at
machine precision.
\end{itemize}

The three constructions are summarised in
Table~\ref{tab:unified_gauge}.  The columns are not independent
discoveries:\ each is a specialisation of the gauge-networks
framework of Marcolli--van~Suijlekom~\cite{marcolli_vs2014}, with
the Perez-Sanchez correction~\cite{perez_sanchez2024,%
perez_sanchez2025} that the continuum limit of such gauge networks
is Wilson Yang--Mills \emph{without} a Higgs sector.

\begin{table}[h]
\centering
\small
\begin{tabular}{|l|c|l|l|c|}
\hline
\textbf{Construction} & $G$ & \textbf{Host manifold} &
\textbf{Vertex Hilbert space} $\mathcal{H}_v$ &
$\frac{1}{4 N_c}$ \\
\hline
Paper~25 & $U(1)$ & $S^3$ Hopf base & $\mathbb{C}$ (scalar) &
$\frac{1}{4}$ \\
Paper~30 & $\mathrm{SU}(2)$ & $S^3 = \mathrm{SU}(2)$ &
$\mathbb{C}^2$ (Dirac) & $\frac{1}{8}$ \\
ST-SU3 & $\mathrm{SU}(3)$ & $S^5$ Bargmann--Segal &
$\bigoplus_N \mathbb{C}^{\dim(N,0)}$ & $\frac{1}{12}$ \\
\hline
\end{tabular}
\caption{The three Wilson lattice gauge constructions as instances of
the Marcolli--van~Suijlekom gauge-network framework on GeoVac
sub-manifolds.  The kinetic coefficient
$\tfrac{1}{4 N_c}$ in the rightmost column is the weak-coupling
expansion coefficient of the Wilson action $S_W = \beta \sum_P \bigl(
1 - \tfrac{1}{N_c} \mathrm{Re}\,\mathrm{tr}\,U_P \bigr)$ around the
trivial vacuum, in the standard normalisation
$\mathrm{tr}(T^a T^b) = \tfrac{1}{2} \delta^{ab}$.  The three rows
verify this coefficient independently:\ Paper~30~\S 5.2 and
Sprint~ST-SU3~\S 4 to machine precision; the abelian Paper~25
case is the trivial $N_c = 1$ limit of either non-abelian
construction.}
\label{tab:unified_gauge}
\end{table}

In each construction the edge Laplacian $L_1 = B^{\top} B$, where
$B$ is the signed node--edge incidence matrix, plays the role of
weak-coupling kinetic operator on $\mathfrak{g}$-valued $1$-cochains.
Paper~25's ``discrete Hodge--1 Laplacian'' is therefore a single
combinatorial object across all three constructions; only its target
Lie algebra changes ($\mathfrak{u}(1)$, $\mathfrak{su}(2)$,
$\mathfrak{su}(3)$).  This is the single mathematical identity
unifying the three sectors at the kinetic level.

\subsubsection*{The two-step SU(3) story}

The $\mathrm{SU}(3)$ row of Table~\ref{tab:unified_gauge} requires a
specific clarification, because an earlier Sprint~5 Track~S5
(April~2026, memo \texttt{debug/s5\_gauge\_structure\_memo.md})
returned a clean negative when asking whether the $S^5$
Bargmann--Segal lattice carries a non-abelian gauge structure
analogous to Paper~25's $U(1)$.  The two findings are not in tension:\
they answer two different questions, and read together they give a
two-step honest picture.

\begin{itemize}\setlength{\itemsep}{2pt}
\item \emph{Step 1 (Sprint~5 Track~S5, negative).}  An
adjacency-preserving $\mathrm{SU}(3)$ action on the $(N, 0)$ tower of
the Bargmann graph fails Wilson's fixed-group-on-every-link
requirement, because the transitions $(N, 0) \to (N+1, 0)$ are
Clebsch--Gordan intertwiners between distinct $\mathrm{SU}(3)$
irreducible representations of growing dimension $\dim(N, 0) =
(N+1)(N+2)/2$, not group elements of a single $\mathrm{SU}(3)$.  The
natural shell-respecting gauge group of the Bargmann graph is
$U(1)$, and this is recorded in Paper~25~\S VII.A as a structural
asymmetry between the $S^3$ and $S^5$ sectors.
\item \emph{Step 2 (Sprint~ST-SU3, positive at the gauge level).}  A
\emph{shell-independent} Wilson lattice gauge theory with link
variables $U_e \in \mathrm{SU}(3)$ in the fundamental representation
(not in any of the $(N, 0)$ irreps) is a structurally different
object that bypasses the Step~1 obstruction.  The links live as
external gauge degrees of freedom on edges; they do not act on
shell labels.  Sprint~ST-SU3 verified gauge-invariance of the
Wilson action under node-local
$U_e \mapsto g_{\mathrm{src}(e)} U_e g_{\mathrm{tgt}(e)}^{\dagger}$
to machine precision and computed the kinetic coefficient $1/12$
both symbolically and numerically.
\item \emph{Matter coupling (Sprint~ST-SU3 Theorem~3, mixed).}  When
one tries to couple this Wilson $\mathrm{SU}(3)$ to the natural
shell matter on the $(N, 0)$ tower, the original Step~1 obstruction
is rediscovered:\ a fundamental-representation link cannot uniformly
couple to representations of growing dimension, and the matter
coupling reverts to the Clebsch--Gordan intertwiner structure.
Wilson $\mathrm{SU}(3)$ on the Bargmann graph is therefore
self-consistent as a \emph{pure-gauge} theory but does not admit a
natural matter coupling in the sense of Paper~30's spin-$1/2$
Dirac fermions on the Coulomb graph.
\end{itemize}

The two-step picture is consistent with Paper~24's Coulomb/HO
asymmetry and sharpens it from three layers to four:\ to the prior
asymmetries (spectrum-computing role of $L_0$; calibration~$\pi$;
Wilson gauge-with-natural-matter) one now adds a fourth layer at
which the $S^3$ Coulomb side is special:\ the gauge group acts in
the same representation as the matter, and matter--gauge coupling
is structurally available.  On the $S^5$ Bargmann side the
combinatorial Wilson--Hodge vocabulary is universal but the
gauge--matter coupling is not.

\subsubsection*{Marcolli--vS lineage of the unified content}

Each construction in Table~\ref{tab:unified_gauge} is a specific
gauge-networks instance in the sense of Marcolli and
van~Suijlekom~\cite{marcolli_vs2014}, with the Perez-Sanchez
correction~\cite{perez_sanchez2024,perez_sanchez2025} that the
continuum limit of such a gauge network is Wilson Yang--Mills
\emph{without} a Higgs.  The vertex algebra is
$\mathcal{A}_v = \mathbb{C}$ (or $C(V)$ globally) in all three rows;
the vertex Hilbert space and the gauge group differ row-by-row, as
listed in the table.  In each row the gauge-network spectral action
reduces to the standard Wilson action plus a quadratic kinetic term
in the weak-coupling expansion, with coefficient $1/(4 N_c)$.

The lineage placement is identical to that of
Sec.~\ref{sec:lineage}, and each Wilson construction inherits the
same status:\ each is a published-machinery instance with a
specific manifold and gauge group; what is not duplicated in the
literature is the \emph{simultaneous} existence of all three SM
gauge groups as siblings across GeoVac sub-manifolds.

\subsubsection*{What this is not:\ structural distance to a Standard Model}

Table~\ref{tab:unified_gauge} is a catalogue, not a unified theory,
and the rhetoric rule of CLAUDE.md~\S 1.5 forbids treating it as the
latter.  The structural distance from this catalogue to Connes' SM
construction lives in four named gaps.

\begin{enumerate}\setlength{\itemsep}{2pt}
\item[\textbf{G1}] (Three generations of matter.)\ Each
construction in Table~\ref{tab:unified_gauge} carries a single
matter species:\ Paper~30 has one generation of Dirac fermions
(Paper~14 Tier~2~\cite{paper14}); Sprint~ST-SU3 has one tower of
$(N, 0)$ irreps.  Sprint~4H Track~SM-B (April~2026) tested the naive
shell-to-generation map $\sum_f N_c Q_f^2 = 8 = |\lambda_3|$ and
returned a clean negative:\ the per-generation contribution is
charge-universal across the three SM generations whereas every
per-shell $S^3$ invariant varies with $n$, so no shell$\to$generation
map is consistent with charge-universal per-generation structure.
The $8 = 8$ equality was a numerical coincidence with no
representation-theoretic content.  There is no GeoVac mechanism
currently selecting a generation tripling of the Hilbert space.
\item[\textbf{G2}] (Higgs as inner fluctuation of $D$.)\ In
Connes' SM, the Higgs field arises as the off-diagonal block of an
inner fluctuation $D \mapsto D + [a, D]$ for $a \in
\mathcal{A}_F = \mathbb{C} \oplus \mathbb{H} \oplus M_3(\mathbb{C})$
in an almost-commutative algebra.  Per Perez-Sanchez 2024, the
Marcolli--vS gauge-network spectral action gives YM \emph{without}
a Higgs sector, so any GeoVac Higgs would have to come from an
almost-commutative extension $\mathcal{A}_{\mathrm{GV}} \otimes
M_n(\mathbb{C})$ with a specific off-diagonal-block structure
analogous to $\mathbb{C} \oplus \mathbb{H} \oplus M_3(\mathbb{C})$.
None of the three Wilson constructions selects such a structure.
The natural place a Higgs-style inner fluctuation would have to
live is precisely the $a \cdot b \notin \mathcal{O}_{n_{\max}}$
operator-system obstruction of Sec.~\ref{sec:construction}
Remark~\ref{rem:operator_system} (witness pair $M_{2,1,0}^2$,
$14.9\%$ residual against $\mathcal{O}_{n_{\max}=2}$).  Whether
multiplicative inner fluctuations can be defined on a non-algebra
operator system is open in the NCG
literature~\cite{connes_vs2021,hekkelman2024}.
\item[\textbf{G3}] (Hypercharge / electroweak unification.)\ The SM
relation $Q = T_3 + Y_W/2$ requires co-located $U(1)_Y$ and
$\mathrm{SU}(2)_L$ on a common Hilbert space with a left--right
chiral asymmetry.  Paper~25's $U(1)$ and Paper~30's
$\mathrm{SU}(2)$ are co-located on the same $S^3$ graph (this is
the most concrete near-term unification target, see below), but
the Camporesi--Higuchi Dirac (Sec.~\ref{sec:construction}
Definition~\ref{def:D_GV_spectral}) is non-chiral by default.  The
chirality grading introduced in Sprint~TS Track~E2 (R3.5, full
Dirac with both chiralities) is a $\mathbb{Z}_2$-grading of the
spinor bundle, not the SM weak-isospin chirality, and it does not
distinguish left- and right-handed sectors of a charged matter
field.
\item[\textbf{G4}] (Cross-manifold obstruction, the dominant gap.)\
Connes' SM uses a \emph{single} spectral triple
$(\mathcal{A}_{\mathrm{SM}}, \mathcal{H}_{\mathrm{SM}},
D_{\mathrm{SM}})$ on a single base manifold, with all three SM
gauge groups emerging from inner automorphisms of one
almost-commutative algebra.  The GeoVac construction has, at
present, three separate spectral triples on two different graphs
(the $S^3$ Coulomb graph for $U(1)$ and $\mathrm{SU}(2)$; the $S^5$
Bargmann graph for $\mathrm{SU}(3)$).  There is no GeoVac operator
that geometrically unifies the $S^3$ and $S^5$ graphs into a
single spectral triple containing both as commuting factors:\ at
best the Bargmann and Coulomb graphs can be encoded as orthogonal
factors of a tensor-product Hilbert space, but the resulting
object is a direct sum of two Wilson-YM instances rather than a
unification.  This is the most consequential gap.  It is also the
most directly tied to the four-layer Coulomb/HO asymmetry of
Sec.~\ref{sec:coulomb_ho}:\ the absence of cross-manifold
unification is a structural symptom of the same asymmetry that
forbids a calibration-$\pi$ analog and a spectrum-computing $L_0$
on $S^5$.
\end{enumerate}

A footnote on the GUT-embedding direction:\ Sprint~4H Track~SM-C
(April~2026) tested whether $1/40 = \Delta$ in Paper~2's $K = \pi(B
+ F - \Delta)$ formula appears naturally in $\mathrm{SU}(5)$,
$E_8$, or $\mathrm{Spin}(10)$ as a Chern number, $\eta$-invariant,
or Chern--Simons level, and returned a clean negative:\ the dual-%
Coxeter rewrite $1/40 = (3/8)/15 = \sin^2\theta_W^{\mathrm{GUT}} /
\dim(\mathbf{\bar{5}} \oplus \mathbf{10})$ is tautological,
recycling the integers $5$ and $8$ already present in Paper~2.  The
naive shell-to-generation embeddings $\mathrm{shell}_n \to
\mathbf{\bar{5}}$ fail at the $\mathrm{SU}(3) \times \mathrm{SU}(2)$
branching level.  GeoVac therefore has no current contact with the
GUT-embedding lineage of the SM, beyond the tautological numerical
coincidences.

\subsubsection*{The most concrete near-term unification target}

Of the four gaps, G4 (cross-manifold) is the most consequential and
G3 (electroweak unification) is the most reachable.  Paper~25's
$U(1)$ and Paper~30's $\mathrm{SU}(2)$ already live on the same
underlying graph (the Fock-projected $S^3$ Coulomb graph) and
share the same Camporesi--Higuchi Dirac operator, the same
incidence structure $B$, and the same edge Laplacian $L_1 = B^{\top}
B$.  The two constructions differ only at the gauge-group level,
and Paper~30~\S 5.1 (Result~1) establishes that Paper~25 is
exactly the maximal-torus reduction of Paper~30.  The natural
unified construction is therefore an $\mathrm{SU}(2) \times U(1)$
Wilson lattice gauge theory on the same graph, with link variables
$U_e \in \mathrm{SU}(2) \times U(1)$ and a chosen embedding of
hypercharge as the $U(1)$ phase.  Whether such a construction
admits a Higgs-style inner fluctuation through an
almost-commutative extension (closing G2 within the electroweak
sector) is an open structural question; whether the
Camporesi--Higuchi Dirac admits a chirality grading compatible with
weak-isospin (closing G3) is the second open structural question.
This is a clean $S^3$-only target that does not require any
progress on the $S^5$ side and that would, if successful, take
GeoVac from ``three Wilson-YM instances on three sub-manifolds''
to ``unified electroweak Wilson construction on $S^3$, plus a
separate $\mathrm{SU}(3)$ instance on $S^5$.''

We do not claim such a unification has been achieved, or even that
its component pieces have been verified.  The status of this
subsection is:\ a recorded structural fact (the unified $1/(4
N_c)$ coefficient and the lineage placement are real); a recorded
honest gap analysis (G1--G4 are real); and a recorded open
direction ($S^3$ electroweak co-location is the most concrete
target).  Per CLAUDE.md~\S 1.5, no claim is made that GeoVac
contains the Standard Model; the claim is only that it contains
the three SM gauge groups as separate Wilson lattice gauge
theories on its own sub-manifolds, with a unified weak-coupling
kinetic coefficient and a unified gauge-network lineage.

% ===================================================================
\section{Marcolli--van~Suijlekom Lineage}
\label{sec:lineage}
% ===================================================================

The construction of Sec.~\ref{sec:construction} is structurally a
specialization of the gauge networks framework of Marcolli and
van~Suijlekom~\cite{marcolli_vs2014}, with the
Perez-Sanchez correction~\cite{perez_sanchez2024,perez_sanchez2025}
that the continuum limit is Yang--Mills without Higgs.  The key
parallels:

\begin{itemize}\setlength{\itemsep}{2pt}
\item \emph{Vertex algebra.}  Marcolli--vS take $\mathcal{A}_V = C(V)$
on a vertex set $V$; we take $\mathcal{A}_{\mathrm{GV}} = \mathbb{C}^{V_{\mathrm{Fock}}}$
on the Fock-projected $S^3$ vertex set (Sec.~\ref{sec:construction}).
\item \emph{Edge connection.}  Marcolli--vS take a connection on
edges valued in a finite group~$G$; Paper~25 takes $G = U(1)$ and
Paper~30 takes $G = \mathrm{SU}(2)$.
\item \emph{Spectral action.}  Marcolli--vS show that the spectral
action of the resulting almost-commutative spectral triple equals a
Wilson lattice gauge action on the connection; Paper~30 verifies this
explicitly for $\mathrm{SU}(2)$ at $n_{\max} = 3, 4$.
\item \emph{Continuum limit.}  Per Perez-Sanchez, the continuum limit
of the gauge-network spectral action is Yang--Mills without a Higgs
sector; Paper~28's heat-kernel Seeley--DeWitt expansion is the
spectral-action expansion in this limit.
\end{itemize}

\begin{observation}[Lineage placement]
\label{obs:lineage}
Each structural ingredient of the GeoVac spectral triple has published
precedent: finite almost-commutative spectral triples
(\cite{krajewski1998,paschke_sitarz2000,cacic2009}), graph spectral
triples (\cite{marcolli_vs2014,perez_sanchez2024}), spectral
truncations (\cite{connes_vs2021,hekkelman2022,hekkelman2024}), and
the Standard Model spectral action (\cite{chamseddine_connes2010}).
What is not duplicated in the literature is the $\alpha$ prediction at
$8.8 \times 10^{-8}$ from a discrete spectral construction; that
remains the original empirical claim of Paper~2.
\end{observation}

This is what the working hypothesis WH1 in the project register
upgrades to ``MODERATE--STRONG'' on (CLAUDE.md~\S 1.7): the structural
framing is in the published lineage; the predictive claim is not
duplicated.

% ===================================================================
\section{Open Questions}
\label{sec:open}
% ===================================================================

The construction of Sec.~\ref{sec:construction} and the audit of
Sec.~\ref{sec:axiom_audit} leave the following structural questions
genuinely open.

\paragraph{Q1.  Order one beyond rank 1.}  The order-one condition
holds trivially for the abelian $\mathcal{A}_{\mathrm{GV}}$ and
non-trivially constrains the rank-1 non-abelian extension
$\mathcal{A}_{\mathrm{GV}} \otimes M_2(\mathbb{C})$ of Paper~30.  At
rank~$\ge 2$ (e.g., $M_3(\mathbb{C})$ or sums of matrix algebras as in
the Standard Model triple), order-one becomes a strong constraint and
typically requires modification of the Dirac operator
($\textit{first-order condition} \Rightarrow$ Higgs sector
in~\cite{chamseddine_connes2010}).  Whether GeoVac admits a natural
rank-$\ge 2$ extension---and what its spectral action would predict---is
an open question.  The Bargmann--Segal $S^5$ HO triple of
Paper~24 was probed for an $\mathrm{SU}(3)$ extension in
Sprint~5 (CLAUDE.md~\S 2) with negative outcome; that is one negative
data point, not the full picture.

\paragraph{Q2.  Real structure at finite $n_{\max}$ vs in the limit.}
Proposition~\ref{prop:reality} establishes a real structure on the
continuum spinor bundle that restricts to the truncation.  Whether the
restricted $J_{\mathrm{GV}}^{(n_{\max})}$ exactly satisfies $J^2 =
+\mathbb{1}$ at finite $n_{\max}$ in finite-precision arithmetic, or
only in the continuum limit, is a structural question that the audit
of Sec.~\ref{sec:axiom_audit} did not pin down to machine precision.
A sprint-scale verification at $n_{\max} \le 5$ would close this.

\paragraph{Q3.  KO-dimension and the four-way coincidence.}  Working
hypothesis WH4 asks: is the four-way $S^3$ coincidence (Sec.~\ref{sec:intro})
forced by the $\mathrm{KO}$-dimension~$3$ structure of the present
triple, or is it independent geometric data?  The relevant test is
whether any other rank-1 spectral triple of $\mathrm{KO}$-dimension~3
fails to produce all four roles.  $S^3$ is unique among compact rank-1
non-abelian simply-connected groups, so the question reduces to:
within the Marcolli--vS gauge-network setting, does
$\mathrm{KO}$-dimension~3 + non-abelian fiber + non-trivial Hopf
bundle uniquely select $S^3$?  The lineage literature does not
appear to address this directly.

\paragraph{Q4.  Spectral-action selection of $\Lambda$.}  Paper~28's
spectral-action computation produced an exact two-term form
$\mathrm{Tr}\,e^{-D^2/\Lambda^2} = (\sqrt{\pi}/2)\Lambda^3 -
(\sqrt{\pi}/4)\Lambda + O(e^{-\pi^2/\Lambda^2})$ but did not
autonomously select $\Lambda$.  Setting the trace equal to $K/\pi$
gave a closed-form Cardano solution $\Lambda_\infty = 3.7102\ldots$
(50 dps), which was tautological in the sense that it was determined
by demanding the equality.  Whether there is a structural principle
that selects $\Lambda$ is open; if so, $K/\pi$ would become a derived
quantity rather than a coincidence.  This is the most direct route by
which Paper~2 could be promoted from observation to derivation, and it
is also the one most resistant to incremental progress.

\paragraph{Q5.  Higher-form structure (Breit / rank-2 connections).}
Paper~25~\S VII.B asked whether the Breit interaction (rank-2 in
electron-electron coupling) lifts the $U(1)$ 1-form picture to a
higher-form gauge theory.  In spectral-triple language this is the
question whether the GeoVac triple admits a natural extension to a
rank-2 fiber that produces the Breit retarded integrals.  Paper~30's
$\mathrm{SU}(2)$ result is rank~1 (vector); rank~2 (tensor) is a
distinct extension; the answer is unknown.

% ===================================================================
\section{Conclusion}
\label{sec:conclusion}
% ===================================================================

This paper has constructed an explicit finite spectral triple
$(\mathcal{A}_{\mathrm{GV}}, \mathcal{H}_{\mathrm{GV}},
D_{\mathrm{GV}})$ on the Fock-projected $S^3$ graph, audited it
against the standard Connes axioms (Table~\ref{tab:axiom_audit},
Theorem~\ref{thm:axiom_complete}), and identified Papers~25, 28, 30,
31 as projections of the resulting object
(Table~\ref{tab:subsectors}).  The construction sits inside the
published Marcolli--van~Suijlekom gauge-network framework with the
Perez-Sanchez correction.  The combination rule
$K = \pi(B + F - \Delta)$ is now phrased as a sum across three
structurally distinct sectors of the same triple
(Observation~\ref{obs:three_sector}), without changing its
conjectural status.  The Coulomb/HO asymmetry of Paper~31 is the
statement that the GeoVac triple is one specific real odd
$\mathrm{KO}$-dimension-3 triple, not a generic class
(Observation~\ref{obs:pi_coulomb}).

Three things this paper did not do.  It did not derive $\alpha$;
Paper~2 stays in Observations.  It did not claim the Standard Model;
GeoVac is rank~1, the SM is higher rank.  It did not prove that the
four-way $S^3$ coincidence is forced by $\mathrm{KO}$-dimension~3;
that is open question Q3.

Three things this paper did do.  It made the spectral-triple framing
explicit and auditable.  It placed the framework in a published
mathematical lineage.  It made the open questions about $\alpha$, the
fiber rank, and the four-way coincidence into precise spectral-triple
questions instead of free-floating conjectures.

% ===================================================================
\begin{thebibliography}{99}

\bibitem{connes1995}
A.~Connes, ``Noncommutative geometry and reality,''
\textit{J.~Math.~Phys.}\ \textbf{36}, 6194--6231 (1995).

\bibitem{connes_marcolli2008}
A.~Connes and M.~Marcolli, \textit{Noncommutative Geometry, Quantum
Fields and Motives}, AMS Colloquium Publications~\textbf{55} (2008).

\bibitem{vansuijlekom_book2015}
W.~D.~van~Suijlekom, \textit{Noncommutative Geometry and Particle
Physics}, Mathematical Physics Studies, Springer (2015).

\bibitem{chamseddine_connes2010}
A.~H.~Chamseddine and A.~Connes,
``Spectral action principle,''
\textit{Commun.~Math.~Phys.}\ \textbf{186}, 731--750 (1997);
``Why the Standard Model,''
\textit{J.~Geom.~Phys.}\ \textbf{58}, 38--47 (2008);
``The spectral action principle in noncommutative geometry and the
superconnection,''
\textit{Phys.~Rev.~D}\ \textbf{83}, 045001 (2011).

\bibitem{marcolli_vs2014}
M.~Marcolli and W.~D.~van~Suijlekom,
``Gauge networks in noncommutative geometry,''
\textit{J.~Geom.~Phys.}\ \textbf{75}, 71--91 (2014),
arXiv:1301.3480.

\bibitem{perez_sanchez2024}
C.~I.~Perez-Sanchez,
``Bratteli networks and the Spectral Action on quivers,''
arXiv:2401.03705 (2024).

\bibitem{perez_sanchez2025}
C.~I.~Perez-Sanchez,
``Comment on `Gauge networks in noncommutative geometry','',
arXiv:2508.17338 (2025).

\bibitem{connes_vs2021}
A.~Connes and W.~D.~van~Suijlekom,
``Spectral truncations in noncommutative geometry and operator
systems,''
\textit{Commun.~Math.~Phys.}\ \textbf{383}, 2021--2067 (2021).

\bibitem{hekkelman2022}
E.~Hekkelman,
``Truncated geometries from spectral triples,''
PhD thesis, Radboud University (2022).

\bibitem{hekkelman2024}
E.~Hekkelman,
``Spectral truncations of the sphere and the Connes--Marcolli
spectral triple,''
arXiv:2403.xxxxx (2024).

\bibitem{krajewski1998}
T.~Krajewski,
``Classification of finite spectral triples,''
\textit{J.~Geom.~Phys.}\ \textbf{28}, 1--30 (1998).

\bibitem{paschke_sitarz2000}
M.~Paschke and A.~Sitarz,
``Discrete spectral triples and their symmetries,''
\textit{J.~Math.~Phys.}\ \textbf{39}, 6191--6205 (1998).

\bibitem{cacic2009}
B.~\'Ca\'ci\'c,
``Real structures on almost-commutative spectral triples,''
\textit{Lett.~Math.~Phys.}\ \textbf{100}, 181--202 (2012).

\bibitem{camporesi_higuchi1996}
R.~Camporesi and A.~Higuchi,
``On the eigenfunctions of the Dirac operator on spheres and real
hyperbolic spaces,''
\textit{J.~Geom.~Phys.}\ \textbf{20}, 1--18 (1996).

\bibitem{paper0}
GeoVac Project, ``Geometric packing and the origin of quantum
numbers,'' Paper~0 (2025).

\bibitem{paper2}
GeoVac Project, ``The fine structure constant from Hopf bundle
spectral invariants,'' Paper~2 (2025; observation status, May~2026).

\bibitem{paper7}
GeoVac Project, ``The Dimensionless Vacuum: $S^3$ Conformal
Equivalence and 18 Symbolic Proofs,'' Paper~7 (2025).

\bibitem{paper14}
GeoVac Project, ``Structurally Sparse Qubit Hamiltonians from Natural
Geometry,'' Paper~14 (2025).

\bibitem{paper18}
GeoVac Project, ``Spectral-Geometric Exchange Constants,'' Paper~18
(2026).

\bibitem{paper22}
GeoVac Project, ``Potential-Independent Angular Sparsity Theorem,''
Paper~22 (2026).

\bibitem{paper23}
GeoVac Project, ``Nuclear Shell Model Qubit Hamiltonians and the
Fock Projection Rigidity Theorem,'' Paper~23 (2026).

\bibitem{paper24}
GeoVac Project, ``The Bargmann--Segal Lattice: $\pi$-Free
Discretization of the 3D Harmonic Oscillator on $S^5$,'' Paper~24
(2026).

\bibitem{paper25}
GeoVac Project, ``The Hopf Graph as a Lattice Gauge Structure: A
Framework Observation,'' Paper~25 (2026).

\bibitem{paper28}
GeoVac Project, ``QED on $S^3$: Spectral Geometry of Perturbative
Transcendentals,'' Paper~28 (2026).

\bibitem{paper30}
GeoVac Project, ``SU(2) Wilson Lattice Gauge Structure on the GeoVac
Hopf Graph: The Non-Abelian Sibling of Paper~25,'' Paper~30 (2026).

\bibitem{paper31}
GeoVac Project, ``Universal/Coulomb Partition: The $\mathcal{A}/D$
Split of the GeoVac Spectral Triple,'' Paper~31 (2026).

\bibitem{paper34}
GeoVac Project, ``Projection Taxonomy and Empirical Matches
Catalogue,'' Paper~34 (2026).

\bibitem{paper35}
GeoVac Project, ``Time as Projection:\ Klein--Gordon Spectrum on
$S^3 \times \mathbb{R}$ and the Temporal-Window Mechanism for $\pi$,''
Paper~35 (2026).

\end{thebibliography}

\end{document}
